% =====================================================================
%  Laboratorio 1 - Monolito con patrón MVC y DAO (personapi-dotnet)
%  Arquitectura de Software
%  Compilar en Overleaf con pdfLaTeX
% =====================================================================
\documentclass[12pt,letterpaper]{article}

% ---------------------------------------------------------------------
%  Paquetes
% ---------------------------------------------------------------------
\usepackage[utf8]{inputenc}
\usepackage[T1]{fontenc}
% es-noquoting: sin él, babel convierte < y > en atajos de comillas y los DbSet<...> del
% diagrama de clases DAO detienen la compilación.
\usepackage[spanish,es-tabla,es-nodecimaldot,es-noquoting]{babel}
\usepackage{lmodern}
\usepackage[margin=2.5cm]{geometry}
\usepackage{graphicx}
\usepackage{xcolor}
\usepackage{amssymb}
\usepackage{booktabs}
\usepackage{tabularx}
\usepackage{longtable}
\usepackage{array}
\usepackage{enumitem}
\usepackage{listings}
\usepackage{tikz}
\usetikzlibrary{positioning,arrows.meta,shapes.multipart,shapes.geometric,fit,backgrounds,calc}
\usepackage{float}
\usepackage{caption}
\usepackage{fancyhdr}
\usepackage{titlesec}
\usepackage{csquotes}
\usepackage[hidelinks]{hyperref}
\usepackage{url}

% ---------------------------------------------------------------------
%  DATOS EDITABLES (cambiar aquí y se actualiza todo el documento)
% ---------------------------------------------------------------------
\newcommand{\universidad}{Pontificia Universidad Javeriana}
\newcommand{\facultad}{Facultad de Ingeniería}
\newcommand{\departamento}{Departamento de Ingeniería de Sistemas}
\newcommand{\asignatura}{Arquitectura de Software}
\newcommand{\profesor}{Andrés Armando Sánchez Martín}
\newcommand{\integranteA}{Sara Albarracín}
\newcommand{\integranteB}{Juan Diego Rojas}
\newcommand{\integranteC}{Miguel Sarmiento}
\newcommand{\ciudad}{Bogotá D.C.}
\newcommand{\fecha}{Octubre de 2026}
\newcommand{\repoURL}{https://github.com/USUARIO/personapi-dotnet}
\newcommand{\repoTag}{v1.0.0}

% ---------------------------------------------------------------------
%  Estilos
% ---------------------------------------------------------------------
\definecolor{azulj}{HTML}{2C5697}
\definecolor{grisfondo}{HTML}{F5F6F8}
\definecolor{grisborde}{HTML}{C9CED6}
\definecolor{verdecom}{HTML}{3A7D44}
\definecolor{rojostr}{HTML}{A3302B}
\definecolor{moradokw}{HTML}{6A3FA0}

\hypersetup{colorlinks=true,linkcolor=azulj,urlcolor=azulj,citecolor=azulj}

\titleformat{\section}{\Large\bfseries\color{azulj}}{\thesection.}{0.6em}{}
\titleformat{\subsection}{\large\bfseries\color{azulj!85!black}}{\thesubsection}{0.6em}{}
\titleformat{\subsubsection}{\normalsize\bfseries}{\thesubsubsection}{0.6em}{}

\pagestyle{fancy}
\fancyhf{}
\fancyhead[L]{\small \asignatura}
\fancyhead[R]{\small Laboratorio 1 -- personapi-dotnet}
\fancyfoot[C]{\thepage}
\renewcommand{\headrulewidth}{0.4pt}

\setlength{\parskip}{0.5em}
\setlength{\parindent}{0pt}

% Listados de código
\lstdefinestyle{base}{
  basicstyle=\ttfamily\footnotesize,
  backgroundcolor=\color{grisfondo},
  frame=single, rulecolor=\color{grisborde},
  breaklines=true, breakatwhitespace=false,
  showstringspaces=false, tabsize=2,
  numbers=left, numberstyle=\tiny\color{gray}, numbersep=6pt,
  xleftmargin=1.5em, framexleftmargin=1.2em,
  keywordstyle=\color{moradokw}\bfseries,
  commentstyle=\color{verdecom}\itshape,
  stringstyle=\color{rojostr},
  captionpos=b,
  literate={á}{{\'a}}1 {é}{{\'e}}1 {í}{{\'i}}1 {ó}{{\'o}}1 {ú}{{\'u}}1
           {Á}{{\'A}}1 {É}{{\'E}}1 {Í}{{\'I}}1 {Ó}{{\'O}}1 {Ú}{{\'U}}1
           {ñ}{{\~n}}1 {Ñ}{{\~N}}1 {ü}{{\"u}}1
}
\lstdefinelanguage{CSharpX}[Sharp]{C}{
  morekeywords={var,async,await,record,init,get,set,required,nameof,
                Task,IActionResult,ActionResult,IEnumerable,List,DbSet,
                DbContext,builder,app,where,partial,virtual},
}
\lstdefinelanguage{yaml}{
  keywords={true,false,null},
  sensitive=false,
  comment=[l]{\#},
  morestring=[b]', morestring=[b]"
}
\lstdefinelanguage{docker}{
  keywords={FROM,AS,WORKDIR,COPY,RUN,ENV,EXPOSE,ENTRYPOINT,CMD,ARG},
  sensitive=true,
  comment=[l]{\#},
  morestring=[b]"
}
\lstdefinelanguage{json}{
  morestring=[b]", stringstyle=\color{rojostr},
  literate={:}{{{\color{moradokw}{:}}}}1 {,}{{{\color{moradokw}{,}}}}1
}
\lstset{style=base}

% Recuadro para capturas pendientes
\newcommand{\captura}[2][0.85\textwidth]{%
  \begin{tikzpicture}
    \node[draw=grisborde, dashed, fill=grisfondo, minimum width=#1,
          minimum height=5cm, align=center, text=gray]
          {\textbf{Captura de pantalla}\\[2pt]\small #2};
  \end{tikzpicture}%
}

% Imagen con espacio reservado: si el archivo existe se inserta; si no, se dibuja el
% recuadro de \captura con la ruta esperada. Basta con subir la imagen a esa ruta
% (relativa a main.tex) para que aparezca, sin tocar el texto.
%   \imagenOPendiente[ancho]{ruta}{descripción}
\newcommand{\imagenOPendiente}[3][0.85\textwidth]{%
  \IfFileExists{#2}%
    {\includegraphics[width=#1,height=0.6\textheight,keepaspectratio]{#2}}%
    {\captura[#1]{#3\\[2pt]\scriptsize\texttt{\detokenize{#2}}}}%
}

% =====================================================================
\begin{document}
% =====================================================================

% ---------------------------------------------------------------------
%  PORTADA
% ---------------------------------------------------------------------
\begin{titlepage}
  \thispagestyle{empty}
  \centering
  {\large\scshape \universidad\par}
  \vspace{0.2cm}
  {\normalsize \facultad\\ \departamento\par}
  \vspace{0.6cm}
  % \includegraphics[width=0.22\textwidth]{logo.png}\par  % <- Descomentar al subir el logo
  \vspace{2.2cm}
  {\color{azulj}\rule{\linewidth}{1.2pt}\par}
  \vspace{0.4cm}
  {\Huge\bfseries Laboratorio 1\par}
  \vspace{0.3cm}
  {\LARGE Implementación de un Monolito\\ con patrón MVC y DAO\par}
  \vspace{0.3cm}
  {\large\itshape personapi-dotnet\par}
  \vspace{0.2cm}
  {\color{azulj}\rule{\linewidth}{1.2pt}\par}
  \vspace{1.2cm}
  {\normalsize .NET 10 (LTS) \quad$\cdot$\quad ASP.NET Core MVC \quad$\cdot$\quad Entity Framework Core 10\\
   SQL Server 2025 \quad$\cdot$\quad REST \quad$\cdot$\quad Swagger / OpenAPI 3 \quad$\cdot$\quad Docker\par}
  \vfill
  {\large\bfseries Integrantes\par}
  \vspace{0.2cm}
  {\large \integranteA\\ \integranteB\\ \integranteC\par}
  \vspace{1cm}
  {\large\bfseries Profesor\par}
  \vspace{0.2cm}
  {\large \profesor\par}
  \vspace{1cm}
  {\large \asignatura\par}
  {\large \ciudad, \fecha\par}
\end{titlepage}

\tableofcontents
\newpage

% ---------------------------------------------------------------------
%  INTRODUCCIÓN (breve, dentro del marco)
% ---------------------------------------------------------------------
\section{Marco conceptual}

Este laboratorio construye \texttt{personapi-dotnet}, una aplicación monolítica que
expone operaciones CRUD (\emph{Create, Read, Update, Delete}) sobre un modelo de datos de
personas, profesiones, estudios y teléfonos, tanto mediante vistas web (HTML generado en el
servidor) como mediante una API REST documentada con Swagger. A continuación se presentan
los conceptos que sustentan la solución.

\subsection{Arquitectura de software}
La arquitectura de software es la estructura (o conjunto de estructuras) de un sistema,
compuesta por sus elementos, las propiedades visibles de estos y las relaciones entre
ellos \cite{bass}. Las decisiones arquitectónicas determinan atributos de calidad como la
mantenibilidad, la escalabilidad, la desplegabilidad y la facilidad de prueba, y son las
más costosas de cambiar una vez el sistema está construido.

\subsection{Arquitectura monolítica}
Un \textbf{monolito} es una aplicación en la que toda la funcionalidad (presentación,
lógica de negocio y acceso a datos) se compila y se despliega como \emph{una sola unidad}
ejecutable \cite{richardson}. Sus principales características son:

\begin{itemize}[itemsep=2pt]
  \item \textbf{Un único artefacto de despliegue}: en este caso, un ensamblado .NET
        publicado dentro de una imagen Docker.
  \item \textbf{Comunicación interna en proceso}: los componentes se invocan mediante
        llamadas a métodos, no por red, lo que reduce la latencia y la complejidad.
  \item \textbf{Base de datos compartida}: todos los módulos operan sobre el mismo esquema.
\end{itemize}

Un monolito no implica desorden: puede (y debe) estar \emph{estructurado en capas}, de
modo que cada capa tenga una responsabilidad clara y dependa únicamente de la capa
inferior. Esta organización, conocida como \textbf{arquitectura en N capas}
(\emph{N-Tier / N-Layer}), es la recomendada por Microsoft para aplicaciones web no
distribuidas, donde el \emph{middle tier} se divide en presentación, lógica de negocio y
acceso a datos \cite{msarchguide}.

\subsection{Patrón Modelo--Vista--Controlador (MVC)}
MVC es un patrón arquitectónico, propuesto originalmente por Trygve Reenskaug en Smalltalk-80
\cite{reenskaug}, que separa una aplicación interactiva en tres tipos de componentes:

\begin{description}[itemsep=2pt]
  \item[Modelo] Representa los datos y las reglas del dominio. Recupera y almacena el
        estado en un medio persistente (la base de datos).
  \item[Vista] Presenta la información del modelo al usuario. En ASP.NET Core las vistas
        se escriben con el motor \emph{Razor} (archivos \texttt{.cshtml}), que mezcla HTML
        con código C\# precedido de \texttt{@}.
  \item[Controlador] Recibe la petición del usuario, interactúa con el modelo y selecciona
        la vista (o la respuesta JSON) que se devolverá.
\end{description}

ASP.NET Core MVC implementa el patrón \emph{Front Controller}: todas las peticiones HTTP
entran por un único punto y el motor de enrutamiento las asigna a una clase controladora
y a un método de acción según la URL (por convención \texttt{/\{controller\}/\{action\}/\{id?\}}
o mediante \emph{attribute routing} con \texttt{[Route]}). Las operaciones de escritura
siguen un esquema de dos pasos: una acción \texttt{HttpGet} presenta el formulario y una
acción \texttt{HttpPost} recibe los datos, que el \emph{model binder} convierte
automáticamente en parámetros u objetos fuertemente tipados \cite{aspnetmvc}.

Entre sus ventajas destacan la separación de responsabilidades (\emph{separation of
concerns}), el soporte nativo de inyección de dependencias e inversión de control, la
facilidad para aplicar pruebas automatizadas y el control total sobre el HTML generado.

\subsection{Patrón Data Access Object (DAO) y Repository}
El patrón \textbf{DAO} abstrae y encapsula todo el acceso a la fuente de datos detrás de
una interfaz \cite{corej2ee}. El resto de la aplicación conoce únicamente la interfaz
(por ejemplo \texttt{IPersonaRepository}) y no sabe si los datos provienen de SQL Server,
de un archivo o de un servicio externo. Sus participantes son:

\begin{itemize}[itemsep=2pt]
  \item \textbf{Cliente (BusinessObject)}: el servicio o controlador que necesita datos.
  \item \textbf{Interfaz DAO}: contrato con las operaciones disponibles
        (\texttt{GetAll}, \texttt{GetById}, \texttt{Add}, \texttt{Update}, \texttt{Delete}).
  \item \textbf{Implementación DAO}: clase concreta que usa la tecnología de persistencia
        (aquí, Entity Framework Core).
  \item \textbf{Objeto de transferencia / entidad}: el objeto que viaja entre capas.
\end{itemize}

En el ecosistema .NET este patrón suele denominarse \textbf{Repository} \cite{fowler}; en
este documento ambos términos se usan de forma equivalente, pues la intención es la misma:
desacoplar la lógica de la aplicación de los detalles de persistencia. Al combinarlo con
inyección de dependencias, la implementación concreta se registra una sola vez en
\texttt{Program.cs} y puede sustituirse (por ejemplo, por un \emph{mock} en pruebas) sin
modificar los controladores.

\subsection{Plataforma .NET}
.NET es una plataforma de desarrollo gratuita, de código abierto y multiplataforma
(Windows, Linux y macOS) que sucede al .NET Framework, exclusivo de Windows. Su núcleo es
el \textbf{Common Language Runtime (CLR)}, implementación del estándar \emph{Common
Language Infrastructure} (CLI), que ejecuta el código intermedio (CIL) generado por los
compiladores y provee servicios de seguridad, recolección de basura y manejo de
excepciones; el código que se ejecuta dentro del CLR se denomina \emph{managed code}.

Microsoft publica una versión mayor cada noviembre; las versiones pares son de soporte a
largo plazo (LTS, tres años) y las impares de soporte estándar (STS) \cite{dotnetpolicy}.
Siguiendo la indicación de usar las \emph{últimas versiones estables}, se seleccionó
\textbf{.NET 10 (LTS)}, publicado en noviembre de 2025 y con soporte hasta noviembre de
2028, en lugar de .NET 7, cuyo soporte terminó en mayo de 2024.

\subsection{ASP.NET Core}
ASP.NET Core es el \emph{framework} web de .NET. Una misma aplicación puede contener
controladores MVC que devuelven vistas Razor y controladores de API
(\texttt{[ApiController]}) que devuelven JSON, lo que permite cubrir en un solo
monolito tanto la interfaz web como la API REST exigidas en el laboratorio. Incluye un
contenedor de inyección de dependencias, un sistema de configuración por capas
(\texttt{appsettings.json}, variables de entorno) y el servidor web multiplataforma
\emph{Kestrel}.

\subsection{Entity Framework Core (ORM)}
Un \textbf{ORM} (\emph{Object-Relational Mapper}) mapea las tablas de una base de datos
relacional a clases del lenguaje y evita escribir código repetitivo de acceso a datos.
Entity Framework Core es el ORM oficial de .NET \cite{efcore}. Sus elementos principales son:

\begin{itemize}[itemsep=2pt]
  \item \textbf{DbContext}: representa una sesión con la base de datos, implementa los
        patrones \emph{Unit of Work} y \emph{Identity Map} y rastrea los cambios.
  \item \textbf{DbSet<T>}: representa una tabla y permite consultas con LINQ, que EF
        traduce a SQL y cuyos resultados mapea a entidades.
  \item \textbf{Proveedor}: \texttt{Microsoft.EntityFrameworkCore.SqlServer} traduce las
        operaciones al dialecto T-SQL.
\end{itemize}

EF Core admite dos enfoques: \emph{Code First} (las clases generan el esquema) y
\emph{Database First}. En este laboratorio se usa \textbf{Database First} mediante
\emph{scaffolding} (ingeniería inversa): el comando \texttt{Scaffold-DbContext} lee el
esquema existente de \texttt{persona\_db} y genera las clases entidad y el contexto.

\subsection{Microsoft SQL Server}
SQL Server es el sistema gestor de bases de datos relacionales de Microsoft. Se utilizó
\textbf{SQL Server 2025} (versión 17.x), la versión estable más reciente, en su edición
\emph{Express} para la instalación local y mediante la imagen oficial
\texttt{mcr.microsoft.com/mssql/server:2025-latest} para el entorno Docker. Su
administración se realiza con \textbf{SQL Server Management Studio (SSMS)}.

\subsection{REST}
REST (\emph{Representational State Transfer}) es un estilo arquitectónico definido por
Roy Fielding \cite{fielding} para sistemas hipermedia distribuidos. Sus restricciones son:
cliente-servidor, comunicación sin estado, respuestas cacheables, interfaz uniforme,
sistema en capas y, opcionalmente, código bajo demanda. En una API REST sobre HTTP:

\begin{itemize}[itemsep=2pt]
  \item Cada recurso se identifica con una URI (\texttt{/api/personas/123}).
  \item Los verbos HTTP expresan la operación: \texttt{GET} (leer), \texttt{POST} (crear),
        \texttt{PUT} (reemplazar/actualizar) y \texttt{DELETE} (eliminar).
  \item Los códigos de estado comunican el resultado: \texttt{200 OK},
        \texttt{201 Created}, \texttt{204 No Content}, \texttt{400 Bad Request},
        \texttt{404 Not Found}, \texttt{409 Conflict}.
\end{itemize}

\subsection{Swagger y OpenAPI 3}
\textbf{OpenAPI} es una especificación estándar, independiente del lenguaje, para
describir APIs HTTP \cite{openapi}; su versión 3 es la evolución de la antigua
especificación Swagger 2.0. \textbf{Swagger UI} es la herramienta que, a partir de ese
documento, genera una página web interactiva para explorar y probar los
\emph{endpoints}. En .NET se integra con el paquete \texttt{Swashbuckle.AspNetCore}, que
genera el documento OpenAPI 3 en \texttt{/swagger/v1/swagger.json} y publica la interfaz
en \texttt{/swagger}.

\subsection{Docker}
Docker empaqueta una aplicación junto con su entorno de ejecución en una \textbf{imagen}
inmutable, que se ejecuta como \textbf{contenedor} aislado \cite{docker}. \emph{Docker
Compose} permite declarar en un archivo YAML varios servicios relacionados (aquí, la
aplicación y SQL Server) y levantarlos con un solo comando. Esto elimina el problema de
\enquote{funciona en mi máquina} y hace reproducible el despliegue del laboratorio en
cualquier sistema operativo.

\subsection{Git, repositorios y etiquetas}
Git es un sistema de control de versiones distribuido. Un \textbf{tag} es una referencia
inmutable a un \emph{commit} concreto que se utiliza para marcar versiones entregables
(por ejemplo \texttt{\repoTag}); es el mecanismo de entrega solicitado para este
laboratorio.

% ---------------------------------------------------------------------
%  DISEÑO
% ---------------------------------------------------------------------
\section{Diseño}

\subsection{Requerimientos}

\textbf{Funcionales}
\begin{enumerate}[label=RF\arabic*., itemsep=2pt]
  \item Crear, consultar (todos y por llave), actualizar y eliminar \emph{personas}.
  \item Crear, consultar, actualizar y eliminar \emph{profesiones}.
  \item Crear, consultar, actualizar y eliminar \emph{teléfonos}, asociados a una persona.
  \item Crear, consultar, actualizar y eliminar \emph{estudios}, que relacionan una persona
        con una profesión.
  \item Ofrecer las operaciones anteriores mediante vistas web (MVC) y mediante
        \emph{endpoints} REST.
  \item Documentar y permitir probar la API desde Swagger UI.
\end{enumerate}

\textbf{No funcionales}
\begin{enumerate}[label=RNF\arabic*., itemsep=2pt]
  \item Usar las últimas versiones estables: .NET 10 LTS y SQL Server 2025.
  \item Estructurar el monolito en capas con responsabilidades separadas.
  \item Aplicar los patrones MVC y DAO.
  \item Desplegar la solución con Docker y Docker Compose.
  \item Validar los datos de entrada y devolver códigos HTTP adecuados.
\end{enumerate}

\subsection{Modelo de datos}

El modelo, proporcionado originalmente en MySQL, consta de cuatro tablas
(Figura~\ref{fig:er}). La tabla \texttt{estudios} es una entidad asociativa que resuelve
la relación muchos-a-muchos entre \texttt{persona} y \texttt{profesion}, con llave
primaria compuesta (\texttt{id\_prof}, \texttt{cc\_per}). La tabla \texttt{telefono}
mantiene una relación uno-a-muchos con \texttt{persona} a través de la columna
\texttt{duenio}.

\begin{figure}[H]
\centering
\begin{tikzpicture}[
  ent/.style={rectangle split, rectangle split parts=2, draw=azulj, thick,
              rectangle split part fill={azulj!85,white}, font=\ttfamily\small,
              align=left, inner sep=4pt, minimum width=4.2cm},
  rel/.style={thick, draw=black!70}
]
\node[ent] (prof) {\textcolor{white}{\textbf{profesion}}
  \nodepart{two} PK id INT\\ \ \ \ nom VARCHAR(90)\\ \ \ \ des VARCHAR(MAX)};
\node[ent, right=1.6cm of prof] (est) {\textcolor{white}{\textbf{estudios}}
  \nodepart{two} PK,FK id\_prof INT\\ PK,FK cc\_per INT\\ \ \ \ \ \ \ fecha DATE\\ \ \ \ \ \ \ univer VARCHAR(50)};
\node[ent, right=1.6cm of est] (per) {\textcolor{white}{\textbf{persona}}
  \nodepart{two} PK cc INT\\ \ \ \ nombre VARCHAR(45)\\ \ \ \ apellido VARCHAR(45)\\ \ \ \ genero CHAR(1)\\ \ \ \ edad INT};
\node[ent, below=1.2cm of per] (tel) {\textcolor{white}{\textbf{telefono}}
  \nodepart{two} PK num VARCHAR(15)\\ \ \ \ oper VARCHAR(45)\\ FK duenio INT};
\draw[rel] (prof.east) -- node[above,font=\scriptsize]{1} node[below,font=\scriptsize,pos=0.85]{N} (est.west);
\draw[rel] (per.west) -- node[above,font=\scriptsize]{1} node[below,font=\scriptsize,pos=0.85]{N} (est.east);
\draw[rel,dashed] (per.south) -- node[right,font=\scriptsize,pos=0.15]{1} node[right,font=\scriptsize,pos=0.85]{N} (tel.north);
\end{tikzpicture}
\caption{Modelo entidad-relación de \texttt{persona\_db} (adaptado a SQL Server).}
\label{fig:er}
\end{figure}

Como el script original está escrito para MySQL, fue necesario traducir los tipos y
restricciones a T-SQL (Tabla~\ref{tab:tipos}).

\begin{table}[H]
\centering
\caption{Adaptación de tipos de MySQL a SQL Server.}
\label{tab:tipos}
\begin{tabularx}{\textwidth}{l l X}
\toprule
\textbf{MySQL} & \textbf{SQL Server} & \textbf{Justificación} \\
\midrule
\texttt{INT(15)}, \texttt{INT(6)}, \texttt{INT(3)} & \texttt{INT} &
  En SQL Server \texttt{INT} no admite ancho de visualización; el rango
  ($\pm 2\,147\,483\,647$) cubre las cédulas colombianas. \\
\texttt{ENUM('M','F')} & \texttt{CHAR(1)} + \texttt{CHECK} &
  SQL Server no tiene \texttt{ENUM}; la restricción \texttt{CHECK (genero IN ('M','F'))}
  garantiza el mismo dominio. \\
\texttt{TEXT} & \texttt{VARCHAR(MAX)} &
  \texttt{TEXT} está obsoleto en SQL Server; \texttt{VARCHAR(MAX)} es su reemplazo. \\
\texttt{ENGINE=InnoDB}, \texttt{CHARSET latin1} & (se omite) &
  Opciones propias de MySQL; SQL Server usa la intercalación de la base de datos. \\
\bottomrule
\end{tabularx}
\end{table}

Las llaves foráneas conservan la semántica del modelo original (\texttt{NO ACTION}): no
se puede eliminar una persona o profesión que tenga estudios o teléfonos asociados. La
aplicación detecta este caso y responde con \texttt{409 Conflict} en la API y con un
mensaje de error en la vista, en lugar de dejar propagar la excepción.

\subsection{Arquitectura de la solución}

La aplicación es un monolito organizado en cuatro capas lógicas dentro de un único
proyecto ASP.NET Core (Figura~\ref{fig:capas}). Cada capa solo depende de la capa
inmediatamente inferior a través de interfaces, y las dependencias se resuelven mediante el
contenedor de inyección de dependencias de ASP.NET Core.

\begin{figure}[H]
\centering
\begin{tikzpicture}[
  capa/.style={draw=azulj, thick, rounded corners=3pt, fill=azulj!8,
               minimum width=12.5cm, minimum height=1.9cm, align=center},
  comp/.style={draw=black!50, rounded corners=2pt, fill=white, font=\small,
               minimum height=0.7cm, inner sep=4pt, align=center},
  flecha/.style={-{Stealth[length=3mm]}, thick, draw=black!70},
  node distance=0.55cm
]
\node (cli) [draw=black!50, rounded corners, fill=black!5, minimum width=12.5cm, minimum height=0.9cm, align=center]
  {\textbf{Clientes}: navegador (vistas HTML) \quad|\quad Swagger UI / Postman (JSON)};

\node (p) [capa, below=of cli] {};
\node[anchor=north west, font=\bfseries\small, text=azulj] at (p.north west) {\ Presentación};
\node[comp] at ($(p.center)+(-3.1,-0.3)$) {Controladores MVC\\ + Vistas Razor};
\node[comp] at ($(p.center)+(3.1,-0.3)$) {Controladores API REST\\ \texttt{[ApiController]}};

\node (s) [capa, below=of p] {};
\node[anchor=north west, font=\bfseries\small, text=azulj] at (s.north west) {\ Lógica de negocio};
\node[comp] at ($(s.center)+(0,-0.3)$) {Servicios: \texttt{IPersonaService}, \texttt{IProfesionService}, \ldots};

\node (d) [capa, below=of s] {};
\node[anchor=north west, font=\bfseries\small, text=azulj] at (d.north west) {\ Acceso a datos (DAO)};
\node[comp] at ($(d.center)+(-3.1,-0.3)$) {Interfaces\\ \texttt{I*Repository}};
\node[comp] at ($(d.center)+(3.1,-0.3)$) {Implementaciones\\ \texttt{*Repository} (EF Core)};

\node (m) [capa, below=of d] {};
\node[anchor=north west, font=\bfseries\small, text=azulj] at (m.north west) {\ Modelo / Persistencia};
\node[comp] at ($(m.center)+(-3.1,-0.3)$) {Entidades\\ \texttt{Persona}, \texttt{Telefono}, \ldots};
\node[comp] at ($(m.center)+(3.1,-0.3)$) {\texttt{PersonaDbContext}};

\node (db) [cylinder, shape border rotate=90, aspect=0.25, draw=azulj, thick, fill=azulj!15,
            minimum width=3.2cm, minimum height=1.4cm, below=of m, align=center]
  {\small SQL Server 2025\\ \small\texttt{persona\_db}};

\draw[flecha] (cli) -- node[right,font=\scriptsize]{HTTP} (p);
\draw[flecha] (p) -- (s);
\draw[flecha] (s) -- (d);
\draw[flecha] (d) -- (m);
\draw[flecha] (m) -- node[right,font=\scriptsize]{ADO.NET / TDS} (db);
\end{tikzpicture}
\caption{Arquitectura en capas del monolito \texttt{personapi-dotnet}.}
\label{fig:capas}
\end{figure}

\begin{description}[itemsep=3pt]
  \item[Presentación] Contiene los controladores MVC (que devuelven vistas Razor) y los
        controladores de API (que devuelven JSON). Ambos tipos de controlador reutilizan
        los mismos servicios, por lo que no se duplica lógica.
  \item[Lógica de negocio] Servicios que aplican las reglas del dominio: verificar que la
        llave no exista antes de crear, que la persona dueña de un teléfono exista, que no
        se elimine un registro con dependencias, etc.
  \item[Acceso a datos (DAO)] Interfaces que definen el contrato de persistencia e
        implementaciones que lo cumplen usando EF Core. Es la única capa que conoce el
        \texttt{DbContext}.
  \item[Modelo / Persistencia] Entidades y contexto generados por \emph{scaffolding} desde
        la base de datos.
\end{description}

\subsection{Modelo C4}

El modelo C4 describe la misma arquitectura en niveles de detalle creciente. Los tres
diagramas se derivan del SAD y del SDD de la wiki y se comprueban contra el código del
repositorio; sus fuentes en texto se versionan en
\texttt{docs/wiki/02-arquitectura/diagramas/} y las imágenes exportadas se guardan junto
a este documento, en \texttt{diagramas/}.

\begin{figure}[H]
\centering
\imagenOPendiente[0.9\textwidth]{diagramas/c4-c1-contexto.png}{Diagrama C4 de nivel 1: contexto}
\caption{C4 nivel 1 (contexto): el usuario, el monolito \texttt{personapi-dotnet} y
SQL Server \texttt{persona\_db} como único almacén.}
\label{fig:c4-contexto}
\end{figure}

\begin{figure}[H]
\centering
\imagenOPendiente[0.9\textwidth]{diagramas/c4-c2-contenedores.png}{Diagrama C4 de nivel 2: contenedores}
\caption{C4 nivel 2 (contenedores): la aplicación ASP.NET Core MVC 10 como único
artefacto desplegable, con sus dos canales (vistas Razor y API REST), y la base
\texttt{persona\_db}.}
\label{fig:c4-contenedores}
\end{figure}

\begin{figure}[H]
\centering
\imagenOPendiente[0.9\textwidth]{diagramas/c4-c3-componentes.png}{Diagrama C4 de nivel 3: componentes}
\caption{C4 nivel 3 (componentes) del contenedor web: controladores MVC y API,
servicios, interfaces DAO e implementaciones EF Core.}
\label{fig:c4-componentes}
\end{figure}

\subsection{Estructura del proyecto}

\begin{lstlisting}[language={},numbers=none,caption={Estructura de carpetas del repositorio.}]
personapi-dotnet/
|-- personapi-dotnet.sln
|-- README.md
|-- Dockerfile
|-- docker-compose.yml
|-- .dockerignore
|-- scripts/
|   |-- 01_ddl.sql                  # Creacion de BD y tablas
|   `-- 02_dml.sql                  # Datos de prueba
`-- personapi-dotnet/
    |-- Program.cs                  # Configuracion, DI, Swagger
    |-- appsettings.json            # Cadena de conexion
    |-- README.md                   # Mapa de capas
    |-- Front/                      # Lo que ve el navegador (Razor server-side)
    |   |-- Views/
    |   |   |-- Shared/_Layout.cshtml
    |   |   |-- Personas/  (Index, Details, Create, Edit, Delete)
    |   |   |-- Profesiones/ ...
    |   |   |-- Telefonos/ ...
    |   |   `-- Estudios/ ...
    |   |-- wwwroot/                #   CSS, JS, imagenes
    |   `-- ViewModels/             #   Contrato de Razor
    `-- Back/                       # Logica: cuatro capas
        |-- 01-Borde/               # CAPA 1: unico sitio que habla HTTP
        |   |-- Controllers/        #   MVC
        |   |   |-- HomeController.cs
        |   |   |-- PersonasController.cs
        |   |   |-- ProfesionesController.cs
        |   |   |-- TelefonosController.cs
        |   |   `-- EstudiosController.cs
        |   |-- ControllersApi/     #   REST /api/*
        |   |   |-- PersonasApiController.cs
        |   |   |-- ProfesionesApiController.cs
        |   |   |-- TelefonosApiController.cs
        |   |   `-- EstudiosApiController.cs
        |   `-- Dtos/               #   Contrato de REST
        |-- 02-Aplicacion/          # CAPA 2: reglas de negocio
        |   `-- Services/
        |       |-- Interfaces/  (IPersonaService.cs, ...)
        |       `-- Implementations/  (PersonaService.cs, ...)
        |-- 03-AccesoDatos/         # CAPA 3: DAO, unico sitio con EF Core
        |   `-- Repository/
        |       |-- Interfaces/  (IPersonaRepository.cs, ...)
        |       `-- Implementations/  (PersonaRepository.cs, ...)
        `-- 04-Persistencia/        # CAPA 4: Database First
            `-- Entities/           #   Generado por Scaffold-DbContext
                |-- PersonaDbContext.cs
                |-- Persona.cs
                |-- Profesion.cs
                |-- Estudio.cs
                `-- Telefono.cs
\end{lstlisting}

\subsection{Diagrama de clases (capa de acceso a datos)}

La Figura~\ref{fig:dao} muestra la aplicación del patrón DAO para la entidad
\texttt{Persona}; las demás entidades siguen exactamente la misma estructura.

\begin{figure}[H]
\centering
\begin{tikzpicture}[
  clase/.style={rectangle split, rectangle split parts=2, draw=black!70, thick,
                rectangle split part fill={azulj!15,white}, font=\ttfamily\scriptsize,
                align=left, inner sep=4pt},
  impl/.style={-{Triangle[open,length=3mm]}, dashed, thick},
  usa/.style={-{Stealth[length=2.5mm]}, thick}
]
\node[clase] (ctrl) {\textbf{PersonasApiController}
  \nodepart{two} - \_service : IPersonaService\\ + GetAll() : IActionResult\\ + Get(cc) : IActionResult\\
                 + Post(persona) : IActionResult\\ + Put(cc, persona) : IActionResult\\ + Delete(cc) : IActionResult};
\node[clase, right=1.2cm of ctrl] (isrv) {\textit{<<interface>>}\\ \textbf{IPersonaService}
  \nodepart{two} + GetAllAsync()\\ + GetByIdAsync(cc)\\ + CreateAsync(p)\\ + UpdateAsync(cc, p)\\ + DeleteAsync(cc)};
\node[clase, below=1.2cm of isrv] (srv) {\textbf{PersonaService}
  \nodepart{two} - \_repo : IPersonaRepository};
\node[clase, below=1.2cm of ctrl] (irepo) {\textit{<<interface>>}\\ \textbf{IPersonaRepository}
  \nodepart{two} + GetAllAsync()\\ + GetByIdAsync(cc)\\ + ExistsAsync(cc)\\ + AddAsync(p)\\ + UpdateAsync(p)\\ + DeleteAsync(p)};
\node[clase, below=1.2cm of irepo] (repo) {\textbf{PersonaRepository}
  \nodepart{two} - \_context : PersonaDbContext};
\node[clase, right=1.2cm of repo] (ctx) {\textbf{PersonaDbContext}
  \nodepart{two} + Personas : DbSet<Persona>\\ + Profesions : DbSet<Profesion>\\ + Estudios : DbSet<Estudio>\\ + Telefonos : DbSet<Telefono>};

\draw[usa] (ctrl) -- (isrv);
\draw[impl] (srv) -- (isrv);
\draw[usa] (srv) -- (irepo);
\draw[impl] (repo) -- (irepo);
\draw[usa] (repo) -- (ctx);
\end{tikzpicture}
\caption{Patrón DAO/Repository aplicado a la entidad \texttt{Persona}.}
\label{fig:dao}
\end{figure}

\subsection{Diseño de la API REST}

Todos los recursos siguen la convención \texttt{/api/\{recurso\}}. La entidad
\texttt{estudios}, por tener llave compuesta, recibe ambos componentes en la ruta.

\begin{longtable}{>{\ttfamily}l >{\ttfamily}l p{4.6cm} >{\ttfamily}l}
\caption{Endpoints REST expuestos por la aplicación.}\label{tab:endpoints}\\
\toprule
\normalfont\textbf{Método} & \normalfont\textbf{Ruta} & \textbf{Descripción} & \normalfont\textbf{Éxito} \\
\midrule
\endfirsthead
\toprule
\normalfont\textbf{Método} & \normalfont\textbf{Ruta} & \textbf{Descripción} & \normalfont\textbf{Éxito} \\
\midrule
\endhead
GET    & /api/personas            & Lista todas las personas       & 200 \\
GET    & /api/personas/\{cc\}     & Obtiene una persona            & 200 / 404 \\
POST   & /api/personas            & Crea una persona               & 201 / 409 \\
PUT    & /api/personas/\{cc\}     & Actualiza una persona          & 204 / 404 \\
DELETE & /api/personas/\{cc\}     & Elimina una persona            & 204 / 409 \\
\midrule
GET    & /api/profesiones         & Lista todas las profesiones    & 200 \\
GET    & /api/profesiones/\{id\}  & Obtiene una profesión          & 200 / 404 \\
POST   & /api/profesiones         & Crea una profesión             & 201 / 409 \\
PUT    & /api/profesiones/\{id\}  & Actualiza una profesión        & 204 / 404 \\
DELETE & /api/profesiones/\{id\}  & Elimina una profesión          & 204 / 409 \\
\midrule
GET    & /api/telefonos           & Lista todos los teléfonos      & 200 \\
GET    & /api/telefonos/\{num\}   & Obtiene un teléfono            & 200 / 404 \\
POST   & /api/telefonos           & Crea un teléfono               & 201 / 400 \\
PUT    & /api/telefonos/\{num\}   & Actualiza un teléfono          & 204 / 404 \\
DELETE & /api/telefonos/\{num\}   & Elimina un teléfono            & 204 / 404 \\
\midrule
GET    & /api/estudios                       & Lista todos los estudios & 200 \\
GET    & /api/estudios/\{idProf\}/\{ccPer\}  & Obtiene un estudio       & 200 / 404 \\
POST   & /api/estudios                       & Crea un estudio          & 201 / 400 \\
PUT    & /api/estudios/\{idProf\}/\{ccPer\}  & Actualiza un estudio     & 204 / 404 \\
DELETE & /api/estudios/\{idProf\}/\{ccPer\}  & Elimina un estudio       & 204 / 404 \\
\bottomrule
\end{longtable}

\subsection{Diseño de las vistas MVC}

Para cada entidad existe un controlador MVC con las cinco acciones estándar de
\emph{scaffolding}, cada una asociada a una vista Razor:

\begin{table}[H]
\centering
\caption{Rutas de las vistas web (ejemplo con \texttt{Personas}).}
\begin{tabular}{>{\ttfamily}l >{\ttfamily}l l}
\toprule
\normalfont\textbf{Ruta} & \normalfont\textbf{Verbo} & \textbf{Vista / acción} \\
\midrule
/Personas                 & GET        & \texttt{Index.cshtml}: tabla con todos los registros \\
/Personas/Details/\{id\}  & GET        & \texttt{Details.cshtml}: detalle de un registro \\
/Personas/Create          & GET, POST  & \texttt{Create.cshtml}: formulario de creación \\
/Personas/Edit/\{id\}     & GET, POST  & \texttt{Edit.cshtml}: formulario de edición \\
/Personas/Delete/\{id\}   & GET, POST  & \texttt{Delete.cshtml}: confirmación de borrado \\
\bottomrule
\end{tabular}
\end{table}

Las listas desplegables para seleccionar la persona dueña de un teléfono o la profesión
de un estudio se construyen con \texttt{SelectList}, de modo que el usuario nunca escribe
manualmente una llave foránea.

\subsection{Diseño del despliegue}

\begin{figure}[H]
\centering
\begin{tikzpicture}[
  cont/.style={draw=azulj, thick, rounded corners=3pt, fill=white, minimum width=4.6cm,
               minimum height=2.2cm, align=center, font=\small},
  flecha/.style={-{Stealth[length=3mm]}, thick}
]
\node[cont] (app) {\textbf{personapi-dotnet}\\ \texttt{aspnet:10.0}\\ puerto 8080};
\node[cont, right=2.2cm of app] (sql) {\textbf{sqlserver}\\ \texttt{mssql/server:2025-latest}\\ puerto 1433};
\node[cont, below=0.9cm of sql, minimum height=1.1cm, fill=grisfondo] (init) {\textbf{db-init}\\ ejecuta \texttt{scripts/*.sql}};
\node[draw, dashed, rounded corners, inner sep=0.45cm, fit=(app)(sql)(init),
      label={[font=\small\bfseries,text=azulj]above:Docker Compose -- red \texttt{personapi-net}}] {};
\node[left=1.3cm of app, align=center, font=\small] (nav) {Navegador\\ \texttt{localhost:8080}};
\draw[flecha] (nav) -- (app);
\draw[flecha] (app) -- node[above,font=\scriptsize]{TDS} (sql);
\draw[flecha] (init) -- (sql);
\end{tikzpicture}
\caption{Vista de despliegue con Docker Compose.}
\label{fig:deploy}
\end{figure}

% ---------------------------------------------------------------------
%  PROCEDIMIENTO
% ---------------------------------------------------------------------
\section{Procedimiento}

El procedimiento sigue la guía del laboratorio. Donde la guía menciona versiones
anteriores (SQL Server 2019, SSMS 18, Visual Studio 2022, .NET 6/7) se usaron las
versiones estables más recientes, según lo solicitado (Tabla~\ref{tab:versiones}).

\begin{table}[H]
\centering
\caption{Versiones de herramientas utilizadas.}
\label{tab:versiones}
\begin{tabularx}{\textwidth}{l l X}
\toprule
\textbf{Herramienta} & \textbf{Guía} & \textbf{Utilizada} \\
\midrule
SDK de .NET              & .NET 6 / 7     & .NET 10 (LTS) \\
Motor de base de datos   & SQL Server 2019 Express & SQL Server 2025 Express / imagen Docker 2025 \\
Cliente de base de datos & SSMS 18        & SSMS 22 \\
IDE                      & Visual Studio 2022 & Visual Studio 2026 Community (requerido para .NET 10) \\
ORM                      & EF Core        & EF Core 10 \\
Documentación de API     & Swagger 3      & Swashbuckle.AspNetCore (OpenAPI 3) \\
Contenedores             & --             & Docker Desktop + Docker Compose v2 \\
\bottomrule
\end{tabularx}
\end{table}

\subsection{Paso 1. Crear el repositorio Git}
Se creó en GitHub el repositorio público \texttt{AS-personapi-dotnet}, inicializado con un
\texttt{README.md} y un \texttt{.gitignore} para Visual Studio. La solución que contiene
se llama \texttt{personapi-dotnet}.

\begin{figure}[H]
\centering
\imagenOPendiente[0.9\textwidth]{figuras/procedimiento/p1-repositorio-github.jpg}{Repositorio \texttt{AS-personapi-dotnet} en GitHub}
\caption{Paso 1: repositorio \texttt{AS-personapi-dotnet} en GitHub, con la rama \texttt{main}
y las carpetas \texttt{docs/}, \texttt{personapi-dotnet/}, \texttt{scripts/} y \texttt{tests/}.}
\label{fig:paso1-repositorio}
\end{figure}

\subsection{Paso 2 y 3. Instalar SQL Server 2025 Express y SSMS}
\begin{enumerate}[itemsep=2pt]
  \item Descargar el instalador de SQL Server 2025 Express desde el sitio de Microsoft y
        elegir el tipo de instalación \textbf{Básica}. La instancia queda registrada como
        \texttt{localhost\textbackslash SQLEXPRESS}.
  \item Habilitar la autenticación mixta (Windows y SQL Server) y asignar contraseña al
        usuario \texttt{sa}, necesario para conectarse desde Docker.
  \item Instalar SQL Server Management Studio (SSMS) y conectarse a la instancia.
\end{enumerate}

\begin{figure}[H]
\centering
\imagenOPendiente[0.9\textwidth]{figuras/procedimiento/p2-ssms-conexion.png}{Conexión a \texttt{localhost\textbackslash SQLEXPRESS} en SSMS}
\caption{Pasos 2 y 3: SSMS 22 conectado a \texttt{localhost\textbackslash SQLEXPRESS}
(SQL Server 17.0) con autenticación de Windows; la barra de estado indica «Conectado» y la
base \texttt{persona\_db}.}
\label{fig:paso2-ssms}
\end{figure}

\subsection{Paso 4 y 5. Crear la base de datos y las tablas (DDL)}
Se creó la base de datos \texttt{persona\_db}, se asignó como propietario al usuario
\texttt{sa} y se crearon las tablas según el modelo, con las adaptaciones de la
Tabla~\ref{tab:tipos}.

\begin{lstlisting}[language=SQL,caption={Script DDL \texttt{scripts/01\_ddl.sql}.}]
IF DB_ID('persona_db') IS NULL
    CREATE DATABASE persona_db;
GO
ALTER AUTHORIZATION ON DATABASE::persona_db TO sa;
GO
USE persona_db;
GO

CREATE TABLE persona (
    cc        INT          NOT NULL,
    nombre    VARCHAR(45)  NOT NULL,
    apellido  VARCHAR(45)  NOT NULL,
    genero    CHAR(1)      NOT NULL,
    edad      INT          NULL,
    CONSTRAINT PK_persona PRIMARY KEY (cc),
    CONSTRAINT CK_persona_genero CHECK (genero IN ('M', 'F'))
);

CREATE TABLE profesion (
    id   INT           NOT NULL,
    nom  VARCHAR(90)   NOT NULL,
    des  VARCHAR(MAX)  NULL,
    CONSTRAINT PK_profesion PRIMARY KEY (id)
);

CREATE TABLE estudios (
    id_prof  INT          NOT NULL,
    cc_per   INT          NOT NULL,
    fecha    DATE         NULL,
    univer   VARCHAR(50)  NULL,
    CONSTRAINT PK_estudios PRIMARY KEY (id_prof, cc_per),
    CONSTRAINT estudio_persona_fk   FOREIGN KEY (cc_per)  REFERENCES persona (cc),
    CONSTRAINT estudio_profesion_fk FOREIGN KEY (id_prof) REFERENCES profesion (id)
);
CREATE INDEX estudio_persona_fk ON estudios (cc_per);

CREATE TABLE telefono (
    num     VARCHAR(15)  NOT NULL,
    oper    VARCHAR(45)  NOT NULL,
    duenio  INT          NOT NULL,
    CONSTRAINT PK_telefono PRIMARY KEY (num),
    CONSTRAINT telefono_persona_fk FOREIGN KEY (duenio) REFERENCES persona (cc)
);
CREATE INDEX telefono_persona_fk ON telefono (duenio);
GO
\end{lstlisting}

\begin{lstlisting}[language=SQL,caption={Script DML \texttt{scripts/02\_dml.sql} (extracto).}]
USE persona_db;
GO
INSERT INTO persona (cc, nombre, apellido, genero, edad) VALUES
 (1010101010, 'Ana',    'Gómez',    'F', 28),
 (1020202020, 'Carlos', 'Pérez',    'M', 35),
 (1030303030, 'Laura',  'Martínez', 'F', 22);

INSERT INTO profesion (id, nom, des) VALUES
 (1, 'Ingeniería de Sistemas', 'Diseño y desarrollo de sistemas de información'),
 (2, 'Medicina',               'Diagnóstico y tratamiento de enfermedades'),
 (3, 'Derecho',                'Asesoría y representación legal');

INSERT INTO estudios (id_prof, cc_per, fecha, univer) VALUES
 (1, 1010101010, '2020-12-15', 'Pontificia Universidad Javeriana'),
 (2, 1020202020, '2015-06-20', 'Universidad Nacional de Colombia'),
 (1, 1030303030, '2024-07-01', 'Universidad de los Andes');

INSERT INTO telefono (num, oper, duenio) VALUES
 ('3001234567', 'Claro',    1010101010),
 ('3109876543', 'Movistar', 1020202020),
 ('3205551234', 'Tigo',     1030303030);
GO
\end{lstlisting}

\begin{figure}[H]
\centering
\imagenOPendiente[0.9\textwidth]{figuras/procedimiento/p5-ssms-tablas.png}{Tablas de \texttt{persona\_db} en el Explorador de objetos de SSMS}
\caption{Paso 5: Explorador de objetos de SSMS con las cuatro tablas de \texttt{persona\_db}:
\texttt{dbo.estudios}, \texttt{dbo.persona}, \texttt{dbo.profesion} y \texttt{dbo.telefono}.}
\label{fig:paso5-tablas}
\end{figure}

\subsection{Paso 6. Instalar Visual Studio Community}
Se instaló Visual Studio 2026 Community con las cargas de trabajo:
\begin{enumerate}[itemsep=1pt]
  \item Desarrollo de ASP.NET y web.
  \item Almacenamiento y procesamiento de datos.
  \item Plantillas de proyecto y elementos de .NET Framework.
  \item Características avanzadas de ASP.NET.
\end{enumerate}

\subsection{Paso 7. Clonar el repositorio}
\begin{lstlisting}[language=bash,numbers=none]
git clone https://github.com/USUARIO/personapi-dotnet.git
cd personapi-dotnet
\end{lstlisting}

\subsection{Paso 8. Desarrollo en Visual Studio}

\subsubsection{8.1 -- 8.4 Creación del proyecto}
Se creó un proyecto con la plantilla \textbf{Aplicación web de ASP.NET Core
(Modelo-Vista-Controlador)}, con nombre \texttt{personapi-dotnet}, \emph{framework}
\textbf{.NET 10.0}, sin autenticación y sin configuración HTTPS. Equivalente por línea de
comandos:

\begin{lstlisting}[language=bash,numbers=none]
dotnet new sln -n personapi-dotnet
dotnet new mvc -n personapi-dotnet -f net10.0 --auth None --no-https
dotnet sln add personapi-dotnet/personapi-dotnet.csproj
\end{lstlisting}

\subsubsection{8.5 -- 8.6 Conexión desde el Explorador de objetos de SQL Server}
Desde \textbf{Ver $\rightarrow$ Explorador de objetos de SQL Server} se agregó la
conexión a \texttt{localhost\textbackslash SQLEXPRESS} y se verificó con \emph{Probar
conexión}.

\subsubsection{8.7 -- 8.8 Instalación de paquetes NuGet}
\begin{lstlisting}[language=bash,numbers=none]
dotnet add package Microsoft.EntityFrameworkCore
dotnet add package Microsoft.EntityFrameworkCore.SqlServer
dotnet add package Microsoft.EntityFrameworkCore.Tools
dotnet add package Microsoft.EntityFrameworkCore.Design
dotnet add package Swashbuckle.AspNetCore
\end{lstlisting}

\subsubsection{8.9 -- 8.10 Generación de entidades (\emph{scaffolding})}
Se creó la carpeta \texttt{Back/04-Persistencia/Entities} y en la Consola del Administrador de paquetes
se ejecutó:

\begin{lstlisting}[language=bash,numbers=none]
Scaffold-DbContext "Server=localhost\SQLEXPRESS;Database=persona_db;Trusted_Connection=True;TrustServerCertificate=true" Microsoft.EntityFrameworkCore.SqlServer -OutputDir Back/04-Persistencia/Entities
\end{lstlisting}

El comando generó \texttt{PersonaDbContext} y las clases \texttt{Persona},
\texttt{Profesion}, \texttt{Estudio} y \texttt{Telefono}, con sus propiedades de
navegación. Luego se movió la cadena de conexión del contexto a
\texttt{appsettings.json}:

\begin{lstlisting}[language=json,caption={Fragmento de \texttt{appsettings.json}.}]
{
  "ConnectionStrings": {
    "PersonaDb": "Server=localhost\\SQLEXPRESS;Database=persona_db;Trusted_Connection=True;TrustServerCertificate=true"
  }
}
\end{lstlisting}

\subsubsection{8.11 Interfaces (contratos DAO)}
\begin{lstlisting}[language=CSharpX,caption={\texttt{Back/03-AccesoDatos/Repository/Interfaces/IPersonaRepository.cs}.}]
public interface IPersonaRepository
{
    Task<IEnumerable<Persona>> GetAllAsync();
    Task<Persona?> GetByIdAsync(int cc);
    Task<bool> ExistsAsync(int cc);
    Task AddAsync(Persona persona);
    Task UpdateAsync(Persona persona);
    Task DeleteAsync(Persona persona);
}
\end{lstlisting}

\subsubsection{8.12 Repositorios (implementación DAO)}
\begin{lstlisting}[language=CSharpX,caption={\texttt{Back/03-AccesoDatos/Repository/Implementations/PersonaRepository.cs}.}]
public class PersonaRepository : IPersonaRepository
{
    private readonly PersonaDbContext _context;

    public PersonaRepository(PersonaDbContext context) => _context = context;

    public async Task<IEnumerable<Persona>> GetAllAsync() =>
        await _context.Personas.AsNoTracking().ToListAsync();

    public async Task<Persona?> GetByIdAsync(int cc) =>
        await _context.Personas.FindAsync(cc);

    public Task<bool> ExistsAsync(int cc) =>
        _context.Personas.AnyAsync(p => p.Cc == cc);

    public async Task AddAsync(Persona persona)
    {
        _context.Personas.Add(persona);
        await _context.SaveChangesAsync();
    }

    public async Task UpdateAsync(Persona persona)
    {
        _context.Personas.Update(persona);
        await _context.SaveChangesAsync();
    }

    public async Task DeleteAsync(Persona persona)
    {
        _context.Personas.Remove(persona);
        await _context.SaveChangesAsync();
    }
}
\end{lstlisting}

\subsubsection{8.13 Controladores}
Los controladores MVC se generaron con \emph{Agregar $\rightarrow$ Nuevo elemento con
scaffolding $\rightarrow$ Controlador de MVC con vistas que usan Entity Framework} y se
modificaron para depender del servicio en lugar del \texttt{DbContext}. Los
controladores de API se escribieron manualmente:

\begin{lstlisting}[language=CSharpX,caption={\texttt{Back/01-Borde/ControllersApi/PersonasApiController.cs}.}]
[ApiController]
[Route("api/personas")]
public class PersonasApiController : ControllerBase
{
    private readonly IPersonaService _service;

    public PersonasApiController(IPersonaService service) => _service = service;

    [HttpGet]
    public async Task<IActionResult> GetAll() => Ok(await _service.GetAllAsync());

    [HttpGet("{cc:int}")]
    public async Task<IActionResult> Get(int cc)
    {
        var persona = await _service.GetByIdAsync(cc);
        return persona is null ? NotFound() : Ok(persona);
    }

    [HttpPost]
    public async Task<IActionResult> Post(Persona persona)
    {
        if (!await _service.CreateAsync(persona))
            return Conflict($"Ya existe una persona con cc {persona.Cc}");
        return CreatedAtAction(nameof(Get), new { cc = persona.Cc }, persona);
    }

    [HttpPut("{cc:int}")]
    public async Task<IActionResult> Put(int cc, Persona persona)
    {
        if (cc != persona.Cc) return BadRequest();
        return await _service.UpdateAsync(cc, persona) ? NoContent() : NotFound();
    }

    [HttpDelete("{cc:int}")]
    public async Task<IActionResult> Delete(int cc)
    {
        var result = await _service.DeleteAsync(cc);
        return result switch
        {
            DeleteResult.NotFound      => NotFound(),
            DeleteResult.HasDependents => Conflict("La persona tiene estudios o teléfonos asociados"),
            _                          => NoContent()
        };
    }
}
\end{lstlisting}

El registro de dependencias y la configuración de Swagger se realizan en
\texttt{Program.cs}:

\begin{lstlisting}[language=CSharpX,caption={Fragmento de \texttt{Program.cs}.}]
var builder = WebApplication.CreateBuilder(args);

builder.Services.AddControllersWithViews()
    .AddJsonOptions(o => o.JsonSerializerOptions.ReferenceHandler =
        System.Text.Json.Serialization.ReferenceHandler.IgnoreCycles);

builder.Services.AddDbContext<PersonaDbContext>(options =>
    options.UseSqlServer(builder.Configuration.GetConnectionString("PersonaDb")));

// Capa de acceso a datos (DAO)
builder.Services.AddScoped<IPersonaRepository, PersonaRepository>();
builder.Services.AddScoped<IProfesionRepository, ProfesionRepository>();
builder.Services.AddScoped<ITelefonoRepository, TelefonoRepository>();
builder.Services.AddScoped<IEstudioRepository, EstudioRepository>();

// Capa de lógica de negocio
builder.Services.AddScoped<IPersonaService, PersonaService>();
builder.Services.AddScoped<IProfesionService, ProfesionService>();
builder.Services.AddScoped<ITelefonoService, TelefonoService>();
builder.Services.AddScoped<IEstudioService, EstudioService>();

// Swagger / OpenAPI 3
builder.Services.AddEndpointsApiExplorer();
builder.Services.AddSwaggerGen(c =>
    c.SwaggerDoc("v1", new() { Title = "personapi-dotnet", Version = "v1" }));

var app = builder.Build();

app.UseSwagger();
app.UseSwaggerUI();
app.UseStaticFiles();
app.UseRouting();
app.MapControllerRoute(
    name: "default",
    pattern: "{controller=Home}/{action=Index}/{id?}");

app.Run();
\end{lstlisting}

\subsubsection{8.14 Despliegue}

El proyecto se despliega de dos formas: como proceso local de ASP.NET Core contra
SQL Server 2025 Express (perfil \texttt{http}, \url{http://localhost:5093}) o con
Docker Compose (\url{http://localhost:8080}). Los dos caminos, con todos sus pasos
desde la instalación de las herramientas, se describen en el manual técnico de la
sección~\ref{sec:manual}.

% ---------------------------------------------------------------------
%  MANUAL TÉCNICO DE DESPLIEGUE (HU-017)
% ---------------------------------------------------------------------
\subsection{Manual técnico de despliegue}
\label{sec:manual}

Este manual permite que cualquier persona, partiendo de un equipo con Windows 10 u 11
de 64 bits sin las herramientas instaladas, despliegue \texttt{personapi-dotnet} y
compruebe que funciona, sin consultar a los autores. Cada paso indica qué se hace, el
comando exacto, el resultado esperado y la captura que lo evidencia.

Hay dos caminos, y ambos parten de los pasos M1 a M7:

\begin{description}[itemsep=2pt]
  \item[Camino A · ejecución local] La aplicación corre con \texttt{dotnet run} contra
        la instancia local \texttt{localhost\textbackslash SQLEXPRESS}, con autenticación
        integrada de Windows. Es el camino con el que se verifican las historias del
        proyecto (pasos M8 a M11).
  \item[Camino B · Docker Compose] Un solo comando levanta SQL Server 2025, carga los
        scripts y publica la aplicación en contenedores. No requiere SQL Server instalado
        en el equipo (pasos M12 a M16).
\end{description}

El \texttt{README.md} del repositorio sigue siendo la guía operativa breve; este manual
la desarrolla paso a paso y no la sustituye. \textbf{Ninguna captura ni comando de este
manual muestra contraseñas}: la del usuario \texttt{sa} del camino B vive solo en el
archivo \texttt{.env}, que Git ignora.

\subsubsection{Requisitos y versiones}

\begin{table}[H]
\centering
\caption{Herramientas del despliegue y versiones con las que se verificó el manual.}
\label{tab:manual-versiones}
\begin{tabularx}{\textwidth}{l l X}
\toprule
\textbf{Herramienta} & \textbf{Camino} & \textbf{Versión verificada} \\
\midrule
Git                         & A y B & 2.53 (\texttt{git --version}) \\
SDK de .NET                 & A     & 10.0.402 (\texttt{dotnet --version}); cualquier SDK 10.0.x \\
SQL Server 2025 Express     & A     & 17.0.1000.7, Express Edition (64-bit) \\
SSMS                        & A     & SQL Server Management Studio 22 \\
\texttt{sqlcmd}             & A     & \texttt{sqlcmd} de Microsoft (versión Go) \\
Visual Studio 2026 Community & A (opcional) & Solo para editar y depurar; no es necesario para desplegar \\
Docker Desktop              & B     & Docker Engine 29.8.1 y Docker Compose v5.5.1, motor Linux sobre WSL~2 \\
Imágenes Docker             & B     & \url{mcr.microsoft.com/dotnet/sdk:10.0},
                                      \url{mcr.microsoft.com/dotnet/aspnet:10.0} y
                                      \url{mcr.microsoft.com/mssql/server:2025-latest}
                                      (SQL Server 17.0.5035.5, Express) \\
\bottomrule
\end{tabularx}
\end{table}

Recursos mínimos para el camino B: 4~GB de RAM libres para Docker y unos 5~GB de disco
para las imágenes. La primera construcción descarga las imágenes y tarda varios
minutos; las siguientes, segundos.

\subsubsection{Paso M1. Instalar Git}
\begin{enumerate}[itemsep=2pt]
  \item Descargar el instalador de \url{https://git-scm.com/downloads} y ejecutarlo con
        las opciones por defecto.
  \item Abrir una terminal nueva (PowerShell) y comprobar la instalación:
\end{enumerate}
\begin{lstlisting}[language=bash,numbers=none]
git --version
# git version 2.53.0.windows.2
\end{lstlisting}

\begin{figure}[H]
\centering
\imagenOPendiente{figuras/manual/01-git-version.png}{Terminal con la salida de \texttt{git --version}}
\caption{Paso M1: Git instalado.}
\end{figure}

\paragraph{Instalación por terminal (\texttt{winget}).} El mismo paso se puede hacer desde PowerShell con el gestor de paquetes de Windows, que descarga el instalador oficial, comprueba su hash y lo ejecuta. Si la herramienta ya está instalada, \texttt{winget} lo detecta y cancela la reinstalación, como ocurrió en el equipo de verificación:

\begin{lstlisting}[language=bash,numbers=none]
winget search --id Git.Git --source winget
winget install --id Git.Git -e --source winget --accept-package-agreements --accept-source-agreements
\end{lstlisting}
\begin{figure}[H]
\centering
\imagenOPendiente[0.85\textwidth]{figuras/manual/m1-a-git-search.png}{\texttt{winget search} encuentra el paquete \texttt{Git.Git}}
\caption{Paso M1: \texttt{winget search} encuentra el paquete \texttt{Git.Git}.}
\end{figure}
\begin{figure}[H]
\centering
\imagenOPendiente[0.85\textwidth]{figuras/manual/m1-b-git-install.png}{\texttt{winget install} de Git; Git ya estaba instalado y \texttt{winget} canceló la reinstalación}
\caption{Paso M1: \texttt{winget install} de Git; Git ya estaba instalado y \texttt{winget} canceló la reinstalación.}
\end{figure}

\paragraph{Asistente gráfico de Git.} Quien prefiera el instalador descargado de la web pasa por las páginas siguientes. En todas se deja la opción que viene marcada y se pulsa \emph{Next}; las que importan para este proyecto son la del PATH, que debe quedar en \emph{Git from the command line and also from 3rd-party software} para que PowerShell y Visual Studio encuentren \texttt{git}, y la de finales de línea, \emph{Checkout Windows-style, commit Unix-style line endings}. Las capturas son del instalador 2.56.0.3 que descarga \texttt{winget} en la fecha de verificación; las páginas y sus valores por defecto son los mismos que en la 2.53. Notas sobre el equipo de verificación:
\begin{itemize}[itemsep=2pt]
  \item Como Git ya estaba instalado, el asistente reutiliza la carpeta de la instalación existente y no muestra la página de carpeta de destino; en un equipo limpio aparece antes de la de componentes.
  \item Con Git instalado, la casilla \emph{Only show new options} viene marcada y salta directo a la instalación; para recorrer todas las páginas hay que desmarcarla en la primera.
  \item El recorrido se canceló en la última página, antes de pulsar \emph{Install}, así que la instalación existente no cambió.
\end{itemize}

\begin{figure}[H]
\centering
\begin{minipage}[t]{0.48\textwidth}
\centering
\imagenOPendiente[\linewidth]{figuras/manual/m1-c01-git-licencia.png}{Licencia GNU GPL v2 del instalador de Git}
\end{minipage}\hfill
\begin{minipage}[t]{0.48\textwidth}
\centering
\imagenOPendiente[\linewidth]{figuras/manual/m1-c02-git-componentes.png}{Selección de componentes de Git}
\end{minipage}
\caption{Paso M1, asistente de Git: licencia y componentes.}
\end{figure}
\begin{figure}[H]
\centering
\begin{minipage}[t]{0.48\textwidth}
\centering
\imagenOPendiente[\linewidth]{figuras/manual/m1-c03-git-editor.png}{Editor por defecto de Git (Vim)}
\end{minipage}\hfill
\begin{minipage}[t]{0.48\textwidth}
\centering
\imagenOPendiente[\linewidth]{figuras/manual/m1-c04-git-rama.png}{Nombre de la rama inicial (Let Git decide)}
\end{minipage}
\caption{Paso M1, asistente de Git: editor por defecto y nombre de la rama inicial.}
\end{figure}
\begin{figure}[H]
\centering
\begin{minipage}[t]{0.48\textwidth}
\centering
\imagenOPendiente[\linewidth]{figuras/manual/m1-c05-git-path.png}{PATH: Git desde la línea de comandos y software de terceros}
\end{minipage}\hfill
\begin{minipage}[t]{0.48\textwidth}
\centering
\imagenOPendiente[\linewidth]{figuras/manual/m1-c06-git-ssh.png}{Cliente SSH: OpenSSH incluido}
\end{minipage}
\caption{Paso M1, asistente de Git: PATH (opción recomendada) y cliente SSH.}
\end{figure}
\begin{figure}[H]
\centering
\begin{minipage}[t]{0.48\textwidth}
\centering
\imagenOPendiente[\linewidth]{figuras/manual/m1-c07-git-https.png}{HTTPS con Windows Secure Channel}
\end{minipage}\hfill
\begin{minipage}[t]{0.48\textwidth}
\centering
\imagenOPendiente[\linewidth]{figuras/manual/m1-c08-git-finales-linea.png}{Finales de línea: checkout Windows, commit Unix}
\end{minipage}
\caption{Paso M1, asistente de Git: biblioteca HTTPS y conversión de finales de línea.}
\end{figure}
\begin{figure}[H]
\centering
\begin{minipage}[t]{0.48\textwidth}
\centering
\imagenOPendiente[\linewidth]{figuras/manual/m1-c09-git-terminal.png}{Terminal de Git Bash: MinTTY}
\end{minipage}\hfill
\begin{minipage}[t]{0.48\textwidth}
\centering
\imagenOPendiente[\linewidth]{figuras/manual/m1-c10-git-pull.png}{Comportamiento de git pull: Merge}
\end{minipage}
\caption{Paso M1, asistente de Git: terminal de Git Bash y comportamiento de \texttt{git pull}.}
\end{figure}
\begin{figure}[H]
\centering
\begin{minipage}[t]{0.48\textwidth}
\centering
\imagenOPendiente[\linewidth]{figuras/manual/m1-c11-git-credenciales.png}{Gestor de credenciales: Git Credential Manager}
\end{minipage}\hfill
\begin{minipage}[t]{0.48\textwidth}
\centering
\imagenOPendiente[\linewidth]{figuras/manual/m1-c12-git-extras.png}{Opciones extra: caché del sistema de archivos}
\end{minipage}
\caption{Paso M1, asistente de Git: gestor de credenciales y opciones extra.}
\end{figure}
\begin{figure}[H]
\centering
\imagenOPendiente[0.6\textwidth]{figuras/manual/m1-c13-git-archivos-en-uso.png}{Lista de procesos bash que usan archivos de Git}
\caption{Paso M1, asistente de Git: al reinstalar o actualizar, el asistente pide cerrar las terminales de Git Bash abiertas antes de \emph{Install}. En un equipo limpio esta página no aparece.}
\end{figure}

\subsubsection{Paso M2. Instalar el SDK de .NET 10}
\begin{enumerate}[itemsep=2pt]
  \item Descargar el \textbf{SDK} (no solo el \emph{runtime}) de .NET 10 para Windows x64
        desde \url{https://dotnet.microsoft.com/download/dotnet/10.0} e instalarlo.
  \item En una terminal nueva:
\end{enumerate}
\begin{lstlisting}[language=bash,numbers=none]
dotnet --version
# 10.0.402   (cualquier 10.0.x es valido)
dotnet --list-sdks
\end{lstlisting}

\begin{figure}[H]
\centering
\imagenOPendiente{figuras/manual/02-dotnet-version.png}{Terminal con \texttt{dotnet --version} y \texttt{dotnet --list-sdks}}
\caption{Paso M2: SDK de .NET 10 instalado.}
\end{figure}

\paragraph{Instalación por terminal (\texttt{winget}).} El mismo paso se puede hacer desde PowerShell con el gestor de paquetes de Windows, que descarga el instalador oficial, comprueba su hash y lo ejecuta. Si la herramienta ya está instalada, \texttt{winget} lo detecta y cancela la reinstalación, como ocurrió en el equipo de verificación:

\begin{lstlisting}[language=bash,numbers=none]
winget search --id Microsoft.DotNet.SDK --source winget
winget install --id Microsoft.DotNet.SDK.10 -e --source winget --accept-package-agreements --accept-source-agreements
\end{lstlisting}
\begin{figure}[H]
\centering
\imagenOPendiente[0.85\textwidth]{figuras/manual/m2-a-dotnet-search.png}{paquetes del SDK de .NET en \texttt{winget}; se usa \texttt{Microsoft.DotNet.SDK.10}}
\caption{Paso M2: paquetes del SDK de .NET en \texttt{winget}; se usa \texttt{Microsoft.DotNet.SDK.10}.}
\end{figure}
\begin{figure}[H]
\centering
\imagenOPendiente[0.85\textwidth]{figuras/manual/m2-b-dotnet-install.png}{\texttt{winget install} del SDK de .NET 10 y \texttt{dotnet --list-sdks}}
\caption{Paso M2: \texttt{winget install} del SDK de .NET 10 y \texttt{dotnet --list-sdks}.}
\end{figure}

\paragraph{Asistente gráfico.} El instalador del SDK (el mismo que descarga \texttt{winget}, versión 10.0.403 en la fecha de verificación) tiene una sola página: la descripción del SDK y los términos que se aceptan al pulsar \emph{Instalar}. La instalación pide permiso de administrador y no tiene opciones. En el equipo de verificación se cerró con \emph{Cerrar} sin instalar, y \texttt{dotnet --list-sdks} siguió mostrando solo la 10.0.402.

\begin{figure}[H]
\centering
\imagenOPendiente[0.6\textwidth]{figuras/manual/m2-c-dotnet-instalador.png}{Instalador del SDK de .NET 10.0.403 con los botones Instalar y Cerrar}
\caption{Paso M2: asistente del SDK de .NET 10; se cerró sin instalar.}
\end{figure}

\subsubsection{Paso M3. Instalar SQL Server 2025 Express (camino A)}
\begin{enumerate}[itemsep=2pt]
  \item Descargar SQL Server 2025 Express desde
        \url{https://www.microsoft.com/sql-server/sql-server-downloads}.
  \item Ejecutar el instalador y elegir el tipo de instalación \textbf{Básica}. Al
        terminar, el instalador muestra el nombre de la instancia,
        \texttt{SQLEXPRESS}, y la cadena de conexión
        \texttt{localhost\textbackslash SQLEXPRESS}.
  \item No hace falta crear usuarios ni contraseñas: el camino A usa la autenticación
        integrada de Windows con el usuario que inició sesión, que el instalador deja
        como administrador de la instancia.
  \item Abrir una terminal nueva y comprobar que el servicio de la instancia está en
        ejecución y que responde con la versión y la edición esperadas (\texttt{sqlcmd}
        se instala en el paso M4):
\end{enumerate}
\begin{lstlisting}[language=bash,numbers=none]
Get-Service 'MSSQL$SQLEXPRESS' | Format-Table Status, Name, DisplayName -AutoSize
# Running  MSSQL$SQLEXPRESS  SQL Server (SQLEXPRESS)
sqlcmd -S 'localhost\SQLEXPRESS' -E -C -W -Q "SELECT SERVERPROPERTY('ProductVersion') AS version, SERVERPROPERTY('Edition') AS edicion"
# 17.0.1000.7 Express Edition (64-bit)
\end{lstlisting}

\begin{figure}[H]
\centering
\imagenOPendiente{figuras/manual/03-sqlserver-servicio.png}{Servicio \texttt{MSSQL\$SQLEXPRESS} en ejecución y versión de la instancia}
\caption{Paso M3: servicio SQL Server (SQLEXPRESS) en ejecución y versión 17.0.1000.7
Express Edition.}
\end{figure}

\paragraph{Instalación por terminal (\texttt{winget}).} El paquete de la edición Express es \texttt{Microsoft.\allowbreak SQLServer.\allowbreak 2025.\allowbreak Express}. Su instalador pide permisos de administrador (Windows muestra el aviso de Control de cuentas de usuario) y abre el mismo asistente descrito arriba, donde se elige \textbf{Básica}:

\begin{lstlisting}[language=bash,numbers=none]
winget search "SQL Server 2025" --source winget
winget install --id Microsoft.SQLServer.2025.Express -e --source winget --accept-package-agreements --accept-source-agreements
\end{lstlisting}
En el equipo de verificación no se ejecutó \texttt{winget install} para SQL Server: la instancia \texttt{SQLEXPRESS} ya existía y \texttt{winget} no la reconoce con ese identificador, así que el comando habría iniciado una segunda instalación.

\begin{figure}[H]
\centering
\imagenOPendiente[0.85\textwidth]{figuras/manual/m3-a-sql-search.png}{paquetes de SQL Server 2025 en \texttt{winget}; el camino A usa la edición Express}
\caption{Paso M3: paquetes de SQL Server 2025 en \texttt{winget}; el camino A usa la edición Express.}
\end{figure}

\paragraph{Asistente gráfico.} El instalador de SQL Server 2025 Express pide permiso de administrador al abrirse y ofrece tres tipos de instalación. Para este proyecto basta \textbf{Basic}, que instala el motor con la configuración por defecto y crea la instancia \texttt{SQLEXPRESS}. Si el idioma de Windows y el formato regional no coinciden, como en el equipo de verificación (es-MX y es-ES), avisa que continuará en inglés: se responde \emph{Yes}. Después vienen la licencia (\emph{Accept}) y la carpeta de instalación, con el espacio necesario y el tamaño de la descarga; \emph{Install} descarga unos 730~MB e instala. En el equipo de verificación el recorrido se cerró en esa última página, antes de \emph{Install}, porque la instancia ya existía.

\begin{figure}[H]
\centering
\begin{minipage}[t]{0.48\textwidth}
\centering
\imagenOPendiente[\linewidth]{figuras/manual/m3-b1-sql-tipo.png}{Tipos de instalación: Basic, Custom y Download Media}
\end{minipage}\hfill
\begin{minipage}[t]{0.48\textwidth}
\centering
\imagenOPendiente[\linewidth]{figuras/manual/m3-b2-sql-idioma.png}{Aviso de idioma: continuar en inglés}
\end{minipage}
\caption{Paso M3, asistente de SQL Server 2025 Express: tipo de instalación y aviso de idioma.}
\end{figure}
\begin{figure}[H]
\centering
\begin{minipage}[t]{0.48\textwidth}
\centering
\imagenOPendiente[\linewidth]{figuras/manual/m3-b3-sql-licencia.png}{Términos de licencia de SQL Server con el botón Accept}
\end{minipage}\hfill
\begin{minipage}[t]{0.48\textwidth}
\centering
\imagenOPendiente[\linewidth]{figuras/manual/m3-b4-sql-ubicacion.png}{Carpeta de instalación con espacio requerido y botón Install}
\end{minipage}
\caption{Paso M3, asistente de SQL Server 2025 Express: licencia y carpeta de instalación (se cerró antes de \emph{Install}).}
\end{figure}

\subsubsection{Paso M4. Instalar SSMS 22 y \texttt{sqlcmd} (camino A)}
\begin{enumerate}[itemsep=2pt]
  \item Instalar SQL Server Management Studio 22 desde
        \url{https://learn.microsoft.com/ssms/install/install}.
  \item Instalar \texttt{sqlcmd} desde
        \url{https://learn.microsoft.com/sql/tools/sqlcmd/sqlcmd-utility}, si no quedó
        instalado con SQL Server.
  \item Abrir SSMS. La primera vez ofrece iniciar sesión con una cuenta de Microsoft o
        de GitHub; no es necesario: se elige \textbf{Omitir y agregar cuentas más tarde}.
  \item Conectarse con \textbf{Nombre del servidor}
        \texttt{localhost\textbackslash SQLEXPRESS}, \textbf{Autenticación} de Windows y
        la opción \textbf{Confiar en el certificado del servidor} marcada. El mismo
        resultado se obtiene abriendo SSMS desde la terminal:
\end{enumerate}
\begin{lstlisting}[language=bash,numbers=none]
& "C:\Program Files\Microsoft SQL Server Management Studio 22\Release\Common7\IDE\SSMS.exe" -S localhost\SQLEXPRESS -d persona_db -C
\end{lstlisting}
SSMS pregunta si se desea conectar al servidor indicado con autenticación de Windows y,
al aceptar, el Explorador de objetos muestra la instancia con sus carpetas. La conexión
también se comprueba con \texttt{sqlcmd}, como en el paso M3; si la terminal no reconoce
\texttt{sqlcmd}, se cierra y se abre una nueva para que cargue el \texttt{PATH}
actualizado por el instalador.

\begin{figure}[H]
\centering
\imagenOPendiente{figuras/manual/04-ssms-conexion.png}{SSMS conectado a \texttt{localhost\textbackslash SQLEXPRESS}}
\caption{Paso M4: SSMS 22 conectado a \texttt{localhost\textbackslash SQLEXPRESS}; el
Explorador de objetos muestra la instancia.}
\end{figure}

\paragraph{Instalación por terminal (\texttt{winget}).} El mismo paso se puede hacer desde PowerShell con el gestor de paquetes de Windows, que descarga el instalador oficial, comprueba su hash y lo ejecuta. Si la herramienta ya está instalada, \texttt{winget} lo detecta y cancela la reinstalación, como ocurrió en el equipo de verificación:

\begin{lstlisting}[language=bash,numbers=none]
winget install --id Microsoft.SQLServerManagementStudio.22 -e --source winget --accept-package-agreements --accept-source-agreements
winget install --id Microsoft.Sqlcmd -e --source winget --accept-package-agreements --accept-source-agreements
sqlcmd --version
\end{lstlisting}
\begin{figure}[H]
\centering
\imagenOPendiente[0.85\textwidth]{figuras/manual/m4-a-ssms-install.png}{\texttt{winget install} de SSMS 22}
\caption{Paso M4: \texttt{winget install} de SSMS 22.}
\end{figure}
\begin{figure}[H]
\centering
\imagenOPendiente[0.85\textwidth]{figuras/manual/m4-b-sqlcmd-install.png}{\texttt{winget install} de \texttt{sqlcmd} y su versión}
\caption{Paso M4: \texttt{winget install} de \texttt{sqlcmd} y su versión.}
\end{figure}

SSMS 22 se instala y se actualiza con el Visual Studio Installer. Con el botón \textbf{Modificar} se ven sus cargas de trabajo; para este laboratorio bastan los componentes principales, que se instalan siempre.

\begin{figure}[H]
\centering
\imagenOPendiente[0.85\textwidth]{figuras/manual/m4-c-ssms-installer.png}{Visual Studio Installer, pantalla Modificar de SQL Server Management Studio 22 (se cerró sin cambios)}
\caption{Paso M4: Visual Studio Installer, pantalla Modificar de SQL Server Management Studio 22 (se cerró sin cambios).}
\end{figure}

\paragraph{Asistentes gráficos.} El instalador de SSMS 22 es un arranque del Visual Studio Installer: pide aceptar la licencia con \emph{Continuar}, descarga y prepara el propio instalador y, si SSMS no está instalado, abre la selección de componentes con \emph{Instalar}. En el equipo de verificación, con SSMS 22.11.1 instalado, ofreció la actualización a la 22.11.2, que se cerró sin aceptar; SSMS quedó en la misma versión. \texttt{sqlcmd} se instala con un paquete MSI: con la herramienta ya instalada, el asistente ofrece \emph{Change}, \emph{Repair} o \emph{Remove}; se canceló y el propio asistente confirma que el sistema no se modificó.

\begin{figure}[H]
\centering
\begin{minipage}[t]{0.48\textwidth}
\centering
\imagenOPendiente[\linewidth]{figuras/manual/m4-d1-ssms-continuar.png}{Visual Studio Installer: antes de empezar, botón Continuar}
\end{minipage}\hfill
\begin{minipage}[t]{0.48\textwidth}
\centering
\imagenOPendiente[\linewidth]{figuras/manual/m4-d2-ssms-preparar.png}{Visual Studio Installer: descarga e instalación del instalador}
\end{minipage}
\caption{Paso M4, instalador de SSMS 22: licencia y preparación del Visual Studio Installer.}
\end{figure}
\begin{figure}[H]
\centering
\begin{minipage}[t]{0.48\textwidth}
\centering
\imagenOPendiente[\linewidth]{figuras/manual/m4-d3-ssms-actualizacion.png}{Actualización disponible de SSMS 22 a 22.11.2}
\end{minipage}\hfill
\begin{minipage}[t]{0.48\textwidth}
\centering
\imagenOPendiente[\linewidth]{figuras/manual/m4-d4-ssms-instalados.png}{Productos instalados en el Visual Studio Installer}
\end{minipage}
\caption{Paso M4, instalador de SSMS 22: aviso de actualización (cerrado sin aceptar) y productos instalados.}
\end{figure}
\begin{figure}[H]
\centering
\begin{minipage}[t]{0.48\textwidth}
\centering
\imagenOPendiente[\linewidth]{figuras/manual/m4-e1-sqlcmd-bienvenida.png}{Bienvenida del asistente de Sqlcmd Tools}
\end{minipage}\hfill
\begin{minipage}[t]{0.48\textwidth}
\centering
\imagenOPendiente[\linewidth]{figuras/manual/m4-e2-sqlcmd-mantenimiento.png}{Opciones Change, Repair y Remove de Sqlcmd Tools}
\end{minipage}
\caption{Paso M4, asistente MSI de \texttt{sqlcmd}: bienvenida y mantenimiento de una instalación existente.}
\end{figure}
\begin{figure}[H]
\centering
\imagenOPendiente[0.6\textwidth]{figuras/manual/m4-e3-sqlcmd-interrumpido.png}{Asistente de Sqlcmd Tools interrumpido sin modificar el sistema}
\caption{Paso M4, asistente de \texttt{sqlcmd} cancelado: el sistema no se modificó.}
\end{figure}

\subsubsection{Paso M5. Instalar Visual Studio 2026 Community (opcional)}
Solo es necesario para abrir y depurar la solución. Desde
\url{https://visualstudio.microsoft.com/} se instala con las cargas de trabajo
\textbf{Desarrollo de ASP.NET y web} y \textbf{Almacenamiento y procesamiento de
datos}. Los pasos de despliegue de este manual se hacen desde la terminal y no lo
requieren. La herramienta \texttt{vswhere}, que se instala con Visual Studio, confirma
qué instalación tiene las dos cargas de trabajo:

\begin{lstlisting}[language=bash,numbers=none]
$vswhere = "${env:ProgramFiles(x86)}\Microsoft Visual Studio\Installer\vswhere.exe"
& $vswhere -products * -requires Microsoft.VisualStudio.Workload.NetWeb Microsoft.VisualStudio.Workload.Data -property displayName
# Visual Studio Community 2026
& $vswhere -products * -requires Microsoft.VisualStudio.Workload.NetWeb Microsoft.VisualStudio.Workload.Data -property catalog_productDisplayVersion
# 18.11.1
\end{lstlisting}

\begin{figure}[H]
\centering
\imagenOPendiente{figuras/manual/05-visualstudio-cargas.png}{Terminal con \texttt{vswhere} y la instalación que tiene las dos cargas de trabajo}
\caption{Paso M5: Visual Studio Community 2026 (18.11.1) con las cargas de ASP.NET y web
y de datos.}
\end{figure}

\paragraph{Instalación por terminal (\texttt{winget}).} El mismo paso se puede hacer desde PowerShell con el gestor de paquetes de Windows, que descarga el instalador oficial, comprueba su hash y lo ejecuta. Si la herramienta ya está instalada, \texttt{winget} lo detecta y cancela la reinstalación, como ocurrió en el equipo de verificación:

\begin{lstlisting}[language=bash,numbers=none]
winget install --id Microsoft.VisualStudio.Community -e --source winget --accept-package-agreements --accept-source-agreements
\end{lstlisting}
\begin{figure}[H]
\centering
\imagenOPendiente[0.85\textwidth]{figuras/manual/m5-a-vs-install.png}{\texttt{winget install} de Visual Studio Community}
\caption{Paso M5: \texttt{winget install} de Visual Studio Community.}
\end{figure}

El Visual Studio Installer muestra los productos instalados. Con \textbf{Modificar} en Visual Studio Community 2026 se marcan las dos cargas de trabajo y se confirma con \textbf{Modificar}, o con \textbf{Instalar} en una instalación nueva:

\begin{figure}[H]
\centering
\imagenOPendiente[0.85\textwidth]{figuras/manual/m5-b-vs-installer.png}{Visual Studio Installer con Build Tools, SSMS 22 y Visual Studio Community 2026}
\caption{Paso M5: Visual Studio Installer con Build Tools, SSMS 22 y Visual Studio Community 2026.}
\end{figure}
\begin{figure}[H]
\centering
\imagenOPendiente[0.85\textwidth]{figuras/manual/m5-c-vs-cargas.png}{carga Desarrollo de ASP.NET y web marcada y, en el panel derecho, Almacenamiento y procesamiento de datos (se cerró sin cambios)}
\caption{Paso M5: carga Desarrollo de ASP.NET y web marcada y, en el panel derecho, Almacenamiento y procesamiento de datos (se cerró sin cambios).}
\end{figure}

\paragraph{Asistente gráfico.} El instalador de Visual Studio Community 2026 es el mismo arranque del Visual Studio Installer que el de SSMS: \emph{Continuar}, preparación y, en un equipo sin Visual Studio, la selección de cargas de trabajo de la figura anterior. Con Visual Studio ya instalado ofreció la actualización de características de octubre de 2026 (18.11.2), que se cerró sin aceptar.

\begin{figure}[H]
\centering
\imagenOPendiente[0.85\textwidth]{figuras/manual/m5-d-vs-actualizacion.png}{Actualización disponible de Visual Studio Community 2026 a 18.11.2}
\caption{Paso M5, instalador de Visual Studio Community 2026: aviso de actualización (cerrado sin aceptar).}
\end{figure}

\subsubsection{Paso M6. Instalar Docker Desktop (camino B)}
\begin{enumerate}[itemsep=2pt]
  \item Descargar Docker Desktop para Windows desde
        \url{https://www.docker.com/products/docker-desktop/} e instalarlo con el motor
        \textbf{WSL~2} (opción por defecto). Si Windows lo pide, reiniciar.
  \item Abrir Docker Desktop y esperar a que indique que el motor está en ejecución
        (\emph{Engine running}).
  \item Comprobar desde la terminal:
\end{enumerate}
\begin{lstlisting}[language=bash,numbers=none]
docker --version
# Docker version 29.8.1, build 4a63305
docker compose version
# Docker Compose version v5.5.1
docker info --format 'Motor: {{.ServerVersion}} ({{.OSType}}, {{.OperatingSystem}})'
# Motor: 29.8.1 (linux, Docker Desktop)
\end{lstlisting}

\begin{figure}[H]
\centering
\imagenOPendiente{figuras/manual/06-docker-version.png}{Terminal con \texttt{docker --version}, \texttt{docker compose version} y el motor en ejecución}
\caption{Paso M6: Docker Desktop instalado, con el motor Linux en ejecución.}
\end{figure}

\paragraph{Instalación por terminal (\texttt{winget}).} El mismo paso se puede hacer desde PowerShell con el gestor de paquetes de Windows, que descarga el instalador oficial, comprueba su hash y lo ejecuta. Si la herramienta ya está instalada, \texttt{winget} lo detecta y cancela la reinstalación, como ocurrió en el equipo de verificación:

\begin{lstlisting}[language=bash,numbers=none]
winget install --id Docker.DockerDesktop -e --source winget --accept-package-agreements --accept-source-agreements
\end{lstlisting}
\begin{figure}[H]
\centering
\imagenOPendiente[0.85\textwidth]{figuras/manual/m6-a-docker-install.png}{\texttt{winget install} de Docker Desktop}
\caption{Paso M6: \texttt{winget install} de Docker Desktop.}
\end{figure}

\paragraph{Asistente gráfico.} El instalador de Docker Desktop (unos 620~MB) pide permiso de administrador y, en un equipo limpio, muestra la página de configuración con la opción \emph{Use WSL 2 instead of Hyper-V}, que debe quedar marcada; \emph{OK} instala y al terminar puede pedir cerrar la sesión de Windows. En el equipo de verificación se lanzó la versión 4.97.0 sobre la 4.93.0 instalada y se canceló: el instalador lo confirma y Docker Engine siguió en la 29.8.1, con los contenedores del proyecto en ejecución.

\begin{figure}[H]
\centering
\begin{minipage}[t]{0.48\textwidth}
\centering
\imagenOPendiente[\linewidth]{figuras/manual/m6-b1-docker-inicializando.png}{Instalador de Docker Desktop 4.97.0 inicializando}
\end{minipage}\hfill
\begin{minipage}[t]{0.48\textwidth}
\centering
\imagenOPendiente[\linewidth]{figuras/manual/m6-b2-docker-cancelado.png}{Instalación cancelada: Docker Desktop 4.93.0 sigue instalado}
\end{minipage}
\caption{Paso M6, instalador de Docker Desktop: inicio y cancelación sin cambios.}
\end{figure}

\subsubsection{Paso M7. Clonar el repositorio}
\begin{lstlisting}[language=bash,numbers=none]
git clone <URL-del-repositorio>.git
cd <carpeta-del-clon>
git checkout v1.0.0     # version entregada
\end{lstlisting}
La URL del repositorio público es \url{\repoURL}, la misma que se indica al final del
procedimiento. La raíz del clon contiene \texttt{personapi-dotnet.sln},
\texttt{Dockerfile}, \texttt{docker-compose.yml}, \texttt{.env.example} y las carpetas
\texttt{scripts/}, \texttt{personapi-dotnet/} y \texttt{tests/}. Todos los comandos
siguientes se ejecutan desde esa raíz.

\begin{figure}[H]
\centering
\imagenOPendiente{figuras/manual/07-git-clone.png}{Terminal con el clon terminado y el listado de la raíz del repositorio}
\caption{Paso M7: repositorio clonado y contenido de su raíz. La captura se tomó antes de
publicar la etiqueta, clonando la rama de esta historia.}
\end{figure}

Para confirmar la rama y conocer la estructura del clon:

\begin{lstlisting}[language=bash,numbers=none]
git branch --show-current
git log --oneline -5
git remote -v
Get-ChildItem -Name
Get-ChildItem scripts -Name
Get-ChildItem personapi-dotnet -Name -Exclude bin,obj
\end{lstlisting}
\begin{figure}[H]
\centering
\imagenOPendiente[0.85\textwidth]{figuras/manual/m7-b-git-log.png}{rama, últimos commits y remoto del clon}
\caption{Paso M7: rama, últimos commits y remoto del clon.}
\end{figure}
\begin{figure}[H]
\centering
\imagenOPendiente[0.62\textwidth]{figuras/manual/m7-c-estructura.png}{contenido de la raíz, de \texttt{scripts/} y del proyecto \texttt{personapi-dotnet/}}
\caption{Paso M7: contenido de la raíz, de \texttt{scripts/} y del proyecto \texttt{personapi-dotnet/}.}
\end{figure}

\subsubsection{Camino A · Paso M8. Crear \texttt{persona\_db}}
Se ejecutan los dos scripts en orden: primero el DDL, que crea la base, las tablas, las
llaves y las restricciones, y después el DML, con tres filas de demostración por tabla.
También pueden abrirse y ejecutarse en SSMS, en el mismo orden.

\begin{lstlisting}[language=bash,numbers=none]
sqlcmd -S "localhost\SQLEXPRESS" -E -C -b -i scripts/01_ddl.sql
sqlcmd -S "localhost\SQLEXPRESS" -E -C -b -i scripts/02_dml.sql
sqlcmd -S "localhost\SQLEXPRESS" -E -C -d persona_db -Q "SELECT COUNT(*) FROM persona; SELECT COUNT(*) FROM profesion; SELECT COUNT(*) FROM telefono; SELECT COUNT(*) FROM estudios"
# 3, 3, 3 y 3 filas
\end{lstlisting}

\textbf{Reejecución.} Los dos scripts son idempotentes: si la base y las filas ya
existen no las duplican ni las modifican, y no muestran error. Para volver al estado
inicial exacto se elimina la base (\texttt{DROP DATABASE persona\_db;}) y se ejecutan de
nuevo los dos scripts.

\begin{figure}[H]
\centering
\imagenOPendiente{figuras/manual/08-scripts-sqlcmd.png}{Terminal con la ejecución de \texttt{01\_ddl.sql}, \texttt{02\_dml.sql} y el conteo de filas}
\caption{Paso M8: creación y carga de \texttt{persona\_db}.}
\end{figure}

\begin{figure}[H]
\centering
\imagenOPendiente{figuras/manual/09-ssms-tablas.png}{Explorador de objetos de SSMS con las cuatro tablas de \texttt{persona\_db}}
\caption{Paso M8: las cuatro tablas de \texttt{persona\_db} en el Explorador de objetos
de SSMS y las tres personas cargadas por el DML.}
\end{figure}

El contenido de cada tabla se revisa con una consulta por tabla, en SSMS o con \texttt{sqlcmd}:

\begin{lstlisting}[language=bash,numbers=none]
sqlcmd -S localhost\SQLEXPRESS -E -C -d persona_db -Q "SELECT * FROM persona"
sqlcmd -S localhost\SQLEXPRESS -E -C -d persona_db -Q "SELECT * FROM profesion"
sqlcmd -S localhost\SQLEXPRESS -E -C -d persona_db -Q "SELECT * FROM telefono"
sqlcmd -S localhost\SQLEXPRESS -E -C -d persona_db -Q "SELECT * FROM estudios"
\end{lstlisting}
\begin{figure}[H]
\centering
\imagenOPendiente[0.85\textwidth]{figuras/manual/m8-c-persona.png}{tabla \texttt{persona} con las tres personas de demostración}
\caption{Paso M8: tabla \texttt{persona} con las tres personas de demostración.}
\end{figure}
\begin{figure}[H]
\centering
\imagenOPendiente[0.85\textwidth]{figuras/manual/m8-d-profesion.png}{tabla \texttt{profesion}}
\caption{Paso M8: tabla \texttt{profesion}.}
\end{figure}
\begin{figure}[H]
\centering
\imagenOPendiente[0.85\textwidth]{figuras/manual/m8-e-telefono.png}{tabla \texttt{telefono}; \texttt{duenio} referencia a \texttt{persona.cc}}
\caption{Paso M8: tabla \texttt{telefono}; \texttt{duenio} referencia a \texttt{persona.cc}.}
\end{figure}
\begin{figure}[H]
\centering
\imagenOPendiente[0.85\textwidth]{figuras/manual/m8-f-estudios.png}{tabla \texttt{estudios}, con llave primaria compuesta (\texttt{id\_prof}, \texttt{cc\_per})}
\caption{Paso M8: tabla \texttt{estudios}, con llave primaria compuesta (\texttt{id\_prof}, \texttt{cc\_per}).}
\end{figure}

El DML se puede ejecutar de nuevo sin duplicar filas:

\begin{figure}[H]
\centering
\imagenOPendiente[0.85\textwidth]{figuras/manual/m8-g-reejecucion.png}{segunda ejecución de \texttt{02\_dml.sql}; los conteos siguen en 3 por tabla}
\caption{Paso M8: segunda ejecución de \texttt{02\_dml.sql}; los conteos siguen en 3 por tabla.}
\end{figure}

\subsubsection{Camino A · Paso M9. Configurar la cadena de conexión}
La aplicación lee la cadena de \texttt{ConnectionStrings:PersonaDb}. En
\texttt{appsettings.json} está vacía a propósito, para no versionar datos de conexión;
el valor local se guarda con \emph{User Secrets}, fuera del repositorio:

\begin{lstlisting}[language=bash,numbers=none]
dotnet user-secrets set "ConnectionStrings:PersonaDb" "Server=localhost\SQLEXPRESS;Database=persona_db;Trusted_Connection=True;TrustServerCertificate=True" --project personapi-dotnet/personapi-dotnet.csproj
dotnet user-secrets list --project personapi-dotnet/personapi-dotnet.csproj
\end{lstlisting}

Esta cadena no contiene contraseña porque usa la autenticación integrada de Windows.
Como alternativa equivalente, sin \emph{User Secrets}, se define la variable de entorno
en la terminal antes de ejecutar:

\begin{lstlisting}[language=bash,numbers=none]
$env:ConnectionStrings__PersonaDb = "Server=localhost\SQLEXPRESS;Database=persona_db;Trusted_Connection=True;TrustServerCertificate=True"
\end{lstlisting}

\begin{figure}[H]
\centering
\imagenOPendiente{figuras/manual/10-user-secrets.png}{Terminal con \texttt{dotnet user-secrets list} mostrando la cadena con autenticación integrada}
\caption{Paso M9: cadena de conexión configurada con User Secrets.}
\end{figure}

\subsubsection{Camino A · Paso M10. Restaurar, compilar y probar}
\begin{lstlisting}[language=bash,numbers=none]
dotnet restore
dotnet build
# Compilación realizado correctamente
dotnet test
# Correctas! - Con error: 0, Superado: 56, Omitido: 0, Total: 56
\end{lstlisting}

Las pruebas de integración usan la misma \texttt{persona\_db}: con la cadena definida
se ejecutan contra la base y dejan los datos como los deja \texttt{02\_dml.sql}; sin
ella se omiten.

\begin{figure}[H]
\centering
\imagenOPendiente{figuras/manual/11-build-test.png}{Terminal con \texttt{dotnet build} sin errores y el resumen de \texttt{dotnet test}}
\caption{Paso M10: compilación y pruebas en verde.}
\end{figure}

Los tres comandos también se pueden ejecutar por separado, para ubicar mejor un error:

\begin{figure}[H]
\centering
\imagenOPendiente[0.85\textwidth]{figuras/manual/m10-a-restore.png}{\texttt{dotnet restore}}
\caption{Paso M10: \texttt{dotnet restore}.}
\end{figure}
\begin{figure}[H]
\centering
\imagenOPendiente[0.85\textwidth]{figuras/manual/m10-b-build.png}{\texttt{dotnet build --no-restore} compila la aplicación y el proyecto de pruebas}
\caption{Paso M10: \texttt{dotnet build --no-restore} compila la aplicación y el proyecto de pruebas.}
\end{figure}
\begin{figure}[H]
\centering
\imagenOPendiente[0.85\textwidth]{figuras/manual/m10-c-test.png}{\texttt{dotnet test --no-build}: 56 pruebas superadas, ninguna con error}
\caption{Paso M10: \texttt{dotnet test --no-build}: 56 pruebas superadas, ninguna con error.}
\end{figure}

\subsubsection{Camino A · Paso M11. Ejecutar la aplicación}
\begin{lstlisting}[language=bash,numbers=none]
dotnet run --project personapi-dotnet/personapi-dotnet.csproj
# Now listening on: http://localhost:5093
# Hosting environment: Development
\end{lstlisting}

Con la aplicación en marcha:
\begin{itemize}[itemsep=1pt]
  \item Vistas web: \url{http://localhost:5093}
  \item Swagger UI: \url{http://localhost:5093/swagger}
\end{itemize}
Se detiene con \texttt{Ctrl+C} en la terminal.

\begin{figure}[H]
\centering
\imagenOPendiente{figuras/manual/12-dotnet-run.png}{Terminal con \texttt{dotnet run} escuchando en \texttt{http://localhost:5093}}
\caption{Paso M11: aplicación en ejecución local.}
\end{figure}

\begin{figure}[H]
\centering
\imagenOPendiente{figuras/manual/13-app-local.png}{Navegador en \texttt{http://localhost:5093} con la página de inicio y el menú}
\caption{Paso M11: página de inicio servida localmente.}
\end{figure}

\begin{figure}[H]
\centering
\imagenOPendiente{figuras/manual/14-swagger-local.png}{Swagger UI en \texttt{http://localhost:5093/swagger}}
\caption{Paso M11: Swagger UI en ejecución local.}
\end{figure}

\subsubsection{Camino B · Paso M12. Crear el archivo \texttt{.env}}
Docker Compose lee la contraseña del usuario \texttt{sa} del contenedor de SQL Server
de un archivo \texttt{.env} en la raíz. Se crea a partir de la plantilla y se edita para
dar valor a \texttt{SA\_PASSWORD}:

\begin{lstlisting}[language=bash,numbers=none]
Copy-Item .env.example .env
notepad .env          # escribir SA_PASSWORD=<contrasena-local> y guardar
git status --short    # .env no aparece: Git lo ignora
\end{lstlisting}

La contraseña debe tener al menos 8 caracteres e incluir tres de estos cuatro grupos:
mayúsculas, minúsculas, dígitos y símbolos; si no, SQL Server no arranca. Es una
contraseña local de laboratorio y \textbf{no se captura ni se escribe en ningún
documento}.

\begin{figure}[H]
\centering
\imagenOPendiente[0.85\textwidth]{figuras/manual/m12-a-env-example.png}{plantilla \texttt{.env.example}; \texttt{SA\_PASSWORD} va vacía y se completa solo en \texttt{.env}, que no se captura}
\caption{Paso M12: plantilla \texttt{.env.example}; \texttt{SA\_PASSWORD} va vacía y se completa solo en \texttt{.env}, que no se captura.}
\end{figure}

\subsubsection{Camino B · Paso M13. Construir y levantar los contenedores}
\begin{lstlisting}[language=bash,numbers=none]
docker compose up --build -d
\end{lstlisting}

El comando construye la imagen de la aplicación con el \texttt{Dockerfile} y arranca
tres servicios en la red \texttt{personapi-net}:

\begin{description}[itemsep=2pt]
  \item[\texttt{sqlserver}] SQL Server 2025 Express. Publica el puerto 1433 y guarda los
        datos en el volumen \texttt{sqldata}. Compose espera a que su comprobación de
        salud (\texttt{SELECT 1}) responda.
  \item[\texttt{db-init}] Ejecuta \texttt{scripts/01\_ddl.sql} y
        \texttt{scripts/02\_dml.sql} contra \texttt{sqlserver} y termina. Como los
        scripts son idempotentes, repetir \texttt{docker compose up} no duplica datos.
  \item[\texttt{personapi}] La aplicación. Arranca solo cuando \texttt{db-init} terminó
        bien y recibe la cadena de conexión por la variable
        \texttt{ConnectionStrings\_\_PersonaDb}, con el usuario \texttt{sa} y la
        contraseña de \texttt{.env}.
\end{description}

Salida esperada al final:
\begin{lstlisting}[language=bash,numbers=none]
[+] up 6/6
 Image personapi-personapi       Built
 Network personapi_personapi-net Created
 Volume personapi_sqldata        Created
 Container personapi-sqlserver-1 Healthy
 Container personapi-db-init-1   Exited
 Container personapi-personapi-1 Started
\end{lstlisting}
La primera vez, antes de estas líneas, Docker descarga las imágenes base y compila la
aplicación dentro de la imagen \texttt{sdk:10.0}; eso tarda varios minutos. En las
siguientes ejecuciones los pasos salen como \texttt{CACHED} y el arranque completo toma
unos segundos.

\begin{figure}[H]
\centering
\imagenOPendiente{figuras/manual/15-compose-up.png}{Terminal con la salida final de \texttt{docker compose up --build -d}}
\caption{Paso M13: contenedores construidos y levantados.}
\end{figure}

Los dos archivos que usa este paso están en la raíz del repositorio:

\begin{lstlisting}[language=docker,caption={\texttt{Dockerfile}.}]
FROM mcr.microsoft.com/dotnet/sdk:10.0 AS build
WORKDIR /src
COPY personapi-dotnet/personapi-dotnet.csproj personapi-dotnet/
RUN dotnet restore personapi-dotnet/personapi-dotnet.csproj
COPY personapi-dotnet/ personapi-dotnet/
RUN dotnet publish personapi-dotnet/personapi-dotnet.csproj -c Release -o /app/publish --no-restore

FROM mcr.microsoft.com/dotnet/aspnet:10.0 AS final
WORKDIR /app
COPY --from=build /app/publish .
# La imagen aspnet:10.0 ya escucha en 8080 (ASPNETCORE_HTTP_PORTS); se declara el puerto.
EXPOSE 8080
ENTRYPOINT ["dotnet", "personapi-dotnet.dll"]
\end{lstlisting}

\begin{lstlisting}[language=yaml,caption={\texttt{docker-compose.yml}.}]
name: personapi

services:
  sqlserver:
    image: mcr.microsoft.com/mssql/server:2025-latest
    environment:
      ACCEPT_EULA: "Y"
      MSSQL_SA_PASSWORD: "${SA_PASSWORD:?Defina SA_PASSWORD en .env}"
      MSSQL_PID: "Express"
    ports:
      - "1433:1433"
    volumes:
      - sqldata:/var/opt/mssql
    healthcheck:
      test: ["CMD-SHELL", "/opt/mssql-tools18/bin/sqlcmd -S localhost -U sa -P \"$$MSSQL_SA_PASSWORD\" -C -Q 'SELECT 1' || exit 1"]
      interval: 10s
      timeout: 5s
      retries: 12
      start_period: 20s
    networks:
      - personapi-net

  db-init:
    image: mcr.microsoft.com/mssql/server:2025-latest
    depends_on:
      sqlserver:
        condition: service_healthy
    volumes:
      - ./scripts:/scripts:ro
    environment:
      SA_PASSWORD: "${SA_PASSWORD:?Defina SA_PASSWORD en .env}"
    entrypoint:
      - bash
      - -c
      - |
        for f in /scripts/01_ddl.sql /scripts/02_dml.sql; do
          echo "Ejecutando $$f"
          /opt/mssql-tools18/bin/sqlcmd -S sqlserver -U sa -P "$$SA_PASSWORD" -C -b -i "$$f" || exit 1
        done
    networks:
      - personapi-net

  personapi:
    build: .
    depends_on:
      db-init:
        condition: service_completed_successfully
    environment:
      ConnectionStrings__PersonaDb: "Server=sqlserver,1433;Database=persona_db;User Id=sa;Password=${SA_PASSWORD:?Defina SA_PASSWORD en .env};TrustServerCertificate=True"
    ports:
      - "8080:8080"
    networks:
      - personapi-net

volumes:
  sqldata:

networks:
  personapi-net:
\end{lstlisting}

Dentro de un contenedor Linux no existe la autenticación integrada de Windows, por eso
este camino usa el usuario \texttt{sa}. La contraseña nunca está en el repositorio:
Compose la toma de \texttt{.env} y la sustituye en
\texttt{\$\{SA\_PASSWORD\}} al arrancar.

\subsubsection{Camino B · Paso M14. Comprobar los contenedores}
\begin{lstlisting}[language=bash,numbers=none]
docker compose ps -a
# personapi   running   0.0.0.0:8080->8080/tcp
# sqlserver   running (healthy)   0.0.0.0:1433->1433/tcp
# db-init     exited (0)
docker compose logs db-init
# Ejecutando /scripts/01_ddl.sql
# Ejecutando /scripts/02_dml.sql
docker compose logs personapi
# Now listening on: http://[::]:8080
\end{lstlisting}

\texttt{db-init} debe terminar con código \texttt{0}: es un servicio de un solo uso, no
un fallo.

\begin{figure}[H]
\centering
\imagenOPendiente{figuras/manual/16-compose-ps.png}{Terminal con \texttt{docker compose ps -a} y los logs de \texttt{db-init}}
\caption{Paso M14: estado de los tres servicios.}
\end{figure}

\begin{figure}[H]
\centering
\imagenOPendiente{figuras/manual/17-docker-desktop-contenedores.png}{Docker Desktop con el proyecto \texttt{personapi} y sus contenedores}
\caption{Paso M14: contenedores en Docker Desktop.}
\end{figure}

Los registros y los recursos creados se revisan desde la terminal:

\begin{lstlisting}[language=bash,numbers=none]
docker compose logs db-init
docker compose logs personapi --tail 15
docker image ls personapi-personapi
docker network ls --filter name=personapi
docker volume ls --filter name=personapi
\end{lstlisting}
\begin{figure}[H]
\centering
\imagenOPendiente[0.85\textwidth]{figuras/manual/m14-b-logs-dbinit.png}{\texttt{db-init} ejecutó \texttt{01\_ddl.sql} y luego \texttt{02\_dml.sql}}
\caption{Paso M14: \texttt{db-init} ejecutó \texttt{01\_ddl.sql} y luego \texttt{02\_dml.sql}.}
\end{figure}
\begin{figure}[H]
\centering
\imagenOPendiente[0.85\textwidth]{figuras/manual/m14-c-logs-personapi.png}{\texttt{personapi} escucha en el puerto 8080 del contenedor y consulta \texttt{persona\_db}}
\caption{Paso M14: \texttt{personapi} escucha en el puerto 8080 del contenedor y consulta \texttt{persona\_db}.}
\end{figure}
\begin{figure}[H]
\centering
\imagenOPendiente[0.85\textwidth]{figuras/manual/m14-d-imagenes-red-volumen.png}{imagen \texttt{personapi-personapi}, red \texttt{personapi\_personapi-net} y volumen \texttt{personapi\_sqldata}}
\caption{Paso M14: imagen \texttt{personapi-personapi}, red \texttt{personapi\_personapi-net} y volumen \texttt{personapi\_sqldata}.}
\end{figure}

Lo mismo se ve en Docker Desktop, filtrando por \texttt{personapi} en cada pestaña:

\begin{figure}[H]
\centering
\imagenOPendiente[0.85\textwidth]{figuras/manual/m14-e-docker-imagenes.png}{Docker Desktop, pestaña Images, con la imagen de la aplicación}
\caption{Paso M14: Docker Desktop, pestaña Images, con la imagen de la aplicación.}
\end{figure}
\begin{figure}[H]
\centering
\imagenOPendiente[0.85\textwidth]{figuras/manual/m14-f-docker-volumen.png}{Docker Desktop, pestaña Volumes, con el volumen de la base}
\caption{Paso M14: Docker Desktop, pestaña Volumes, con el volumen de la base.}
\end{figure}
\begin{figure}[H]
\centering
\imagenOPendiente[0.85\textwidth]{figuras/manual/m14-g-docker-logs.png}{Docker Desktop, registros del contenedor \texttt{personapi-personapi-1}}
\caption{Paso M14: Docker Desktop, registros del contenedor \texttt{personapi-personapi-1}.}
\end{figure}

\subsubsection{Camino B · Paso M15. Usar la aplicación}
\begin{itemize}[itemsep=1pt]
  \item Vistas web: \url{http://localhost:8080}
  \item Swagger UI: \url{http://localhost:8080/swagger}
\end{itemize}

\begin{figure}[H]
\centering
\imagenOPendiente{figuras/manual/18-app-docker.png}{Navegador en \texttt{http://localhost:8080} con la página de inicio}
\caption{Paso M15: aplicación servida desde el contenedor.}
\end{figure}

\begin{figure}[H]
\centering
\imagenOPendiente{figuras/manual/19-swagger-docker.png}{Swagger UI (puerto 8080) ejecutando \texttt{GET /api/personas} con respuesta 200}
\caption{Paso M15: Swagger UI desde el contenedor.}
\end{figure}

Con Swagger UI se recorre un CRUD completo sin otra herramienta: en cada operación se pulsa \textbf{Try it out}, se llenan los parámetros o el cuerpo y se pulsa \textbf{Execute}. La secuencia siguiente crea la profesión 40, la modifica, la consulta, la elimina y comprueba que ya no existe, así que la base queda como estaba:

\begin{figure}[H]
\centering
\imagenOPendiente[0.85\textwidth]{figuras/manual/m15-a-swagger-get-lista.png}{\texttt{GET /api/personas} con respuesta 200 y las tres personas}
\caption{Paso M15: \texttt{GET /api/personas} con respuesta 200 y las tres personas.}
\end{figure}
\begin{figure}[H]
\centering
\imagenOPendiente[0.85\textwidth]{figuras/manual/m15-c-swagger-get-cc.png}{\texttt{GET /api/personas/1010101010} con respuesta 200}
\caption{Paso M15: \texttt{GET /api/personas/1010101010} con respuesta 200.}
\end{figure}
\begin{figure}[H]
\centering
\imagenOPendiente[0.85\textwidth]{figuras/manual/m15-d-swagger-post.png}{\texttt{POST /api/profesiones} con la profesión 40; respuesta 201}
\caption{Paso M15: \texttt{POST /api/profesiones} con la profesión 40; respuesta 201.}
\end{figure}
\begin{figure}[H]
\centering
\imagenOPendiente[0.85\textwidth]{figuras/manual/m15-e-swagger-put.png}{\texttt{PUT /api/profesiones/40}; respuesta 204}
\caption{Paso M15: \texttt{PUT /api/profesiones/40}; respuesta 204.}
\end{figure}
\begin{figure}[H]
\centering
\imagenOPendiente[0.85\textwidth]{figuras/manual/m15-f-swagger-get-id.png}{\texttt{GET /api/profesiones/40} devuelve la profesión modificada}
\caption{Paso M15: \texttt{GET /api/profesiones/40} devuelve la profesión modificada.}
\end{figure}
\begin{figure}[H]
\centering
\imagenOPendiente[0.85\textwidth]{figuras/manual/m15-g-swagger-delete.png}{\texttt{DELETE /api/profesiones/40}; respuesta 204}
\caption{Paso M15: \texttt{DELETE /api/profesiones/40}; respuesta 204.}
\end{figure}
\begin{figure}[H]
\centering
\imagenOPendiente[0.85\textwidth]{figuras/manual/m15-h-swagger-404.png}{\texttt{GET /api/profesiones/40} después de eliminarla; respuesta 404 con \texttt{ProblemDetails}}
\caption{Paso M15: \texttt{GET /api/profesiones/40} después de eliminarla; respuesta 404 con \texttt{ProblemDetails}.}
\end{figure}

La interfaz MVC también responde en el 8080, por ejemplo la página de convenciones de errores:

\begin{figure}[H]
\centering
\imagenOPendiente[0.85\textwidth]{figuras/manual/m15-b-convenciones.png}{página Convenciones de errores en \texttt{http://localhost:8080/Home/Convenciones}}
\caption{Paso M15: página Convenciones de errores en \texttt{http://localhost:8080/Home/Convenciones}.}
\end{figure}

\subsubsection{Camino B · Paso M16. Detener y limpiar}
\begin{lstlisting}[language=bash,numbers=none]
docker compose down       # detiene y elimina los contenedores; conserva los datos
docker compose down -v    # ademas elimina el volumen sqldata: la base se pierde
\end{lstlisting}
Tras \texttt{down -v}, el siguiente \texttt{docker compose up -d} vuelve a crear la base
desde los scripts.

\begin{figure}[H]
\centering
\imagenOPendiente[0.85\textwidth]{figuras/manual/m16-a-compose-down.png}{\texttt{docker compose down} elimina contenedores y red; el volumen se conserva}
\caption{Paso M16: \texttt{docker compose down} elimina contenedores y red; el volumen se conserva.}
\end{figure}
\begin{figure}[H]
\centering
\imagenOPendiente[0.85\textwidth]{figuras/manual/m16-b-compose-up-otra-vez.png}{\texttt{docker compose up -d} vuelve a levantar la pila con los datos del volumen}
\caption{Paso M16: \texttt{docker compose up -d} vuelve a levantar la pila con los datos del volumen.}
\end{figure}

\subsubsection{Verificación posterior al despliegue}
Sirve para los dos caminos; solo cambia el puerto: \texttt{5093} en el camino A y
\texttt{8080} en el B. Las comprobaciones se hacen desde el navegador, desde Swagger UI
o con estos comandos de PowerShell:

\begin{lstlisting}[language=bash,numbers=none]
$b = "http://localhost:8080"     # en el camino A: http://localhost:5093
(Invoke-WebRequest "$b/" -UseBasicParsing).StatusCode                     # 200
(Invoke-WebRequest "$b/swagger/index.html" -UseBasicParsing).StatusCode   # 200
Invoke-RestMethod "$b/api/personas" | ForEach-Object { $_ } | Format-Table cc, nombre, apellido
Invoke-WebRequest "$b/api/profesiones" -Method Post -ContentType "application/json" -Body '{"id":40,"nom":"Arquitectura","des":"Prueba de despliegue"}' -UseBasicParsing | Select-Object StatusCode   # 201
(Invoke-WebRequest "$b/api/profesiones/40" -Method Delete -UseBasicParsing).StatusCode       # 204
\end{lstlisting}

\begin{table}[H]
\centering
\caption{Resultado esperado de la verificación posterior al despliegue.}
\label{tab:manual-verificacion}
\begin{tabularx}{\textwidth}{>{\raggedright\arraybackslash}p{4.6cm} >{\raggedright\arraybackslash}p{4.2cm} >{\raggedright\arraybackslash}X}
\toprule
\textbf{Comprobación} & \textbf{Esperado} & \textbf{Qué demuestra} \\
\midrule
\texttt{GET /}                         & 200 & La aplicación MVC responde \\
\texttt{GET /swagger/index.html}       & 200 & Swagger UI publicado \\
\texttt{GET /swagger/v1/swagger.json}  & 200, OpenAPI 3.0.4 & Contrato OpenAPI 3 \\
\texttt{GET /api/personas}             & 200 y las 3 personas del DML & Conexión real a \texttt{persona\_db} \\
\texttt{POST /api/profesiones} (id 40) & 201 con \texttt{Location} & Escritura en la base \\
\texttt{DELETE /api/profesiones/40}    & 204 & La base queda como estaba \\
\bottomrule
\end{tabularx}
\end{table}

\begin{figure}[H]
\centering
\imagenOPendiente[0.85\textwidth]{figuras/manual/v-a-verificacion-estados.png}{códigos de estado de la página de inicio, Swagger UI, el documento OpenAPI y la API}
\caption{Verificación: códigos de estado de la página de inicio, Swagger UI, el documento OpenAPI y la API.}
\end{figure}
\begin{figure}[H]
\centering
\imagenOPendiente[0.85\textwidth]{figuras/manual/v-b-verificacion-datos.png}{datos de \texttt{/api/personas} y \texttt{/api/profesiones} leídos de \texttt{persona\_db}}
\caption{Verificación: datos de \texttt{/api/personas} y \texttt{/api/profesiones} leídos de \texttt{persona\_db}.}
\end{figure}

\subsubsection{Fallos frecuentes}

\begin{longtable}{>{\raggedright\arraybackslash}p{4.4cm} >{\raggedright\arraybackslash}p{4.8cm} >{\raggedright\arraybackslash}p{5.4cm}}
\caption{Fallos frecuentes del despliegue, su causa y su solución.}\label{tab:manual-fallos}\\
\toprule
\textbf{Síntoma} & \textbf{Causa} & \textbf{Solución} \\
\midrule
\endfirsthead
\toprule
\textbf{Síntoma} & \textbf{Causa} & \textbf{Solución} \\
\midrule
\endhead
\texttt{/api/*} responde 500 y el log dice \emph{No se ha inicializado la propiedad
ConnectionString} &
La cadena \texttt{ConnectionStrings:}\allowbreak\texttt{PersonaDb} no está definida
(camino A) &
Repetir el paso M9 en el mismo usuario y proyecto, o definir la variable de entorno en
la misma terminal donde se ejecuta \texttt{dotnet run}. \\
\midrule
\texttt{/api/*} responde 500 y el log dice \emph{Cannot open database} con el nombre \texttt{persona\_db} &
La base no existe o la cadena apunta a otra base &
Ejecutar el paso M8 y revisar el nombre \texttt{Database=persona\_db} de la cadena. \\
\midrule
\texttt{sqlcmd} o la aplicación no conectan: \emph{A network-related or
instance-specific error} &
Conexión rechazada: el servicio \texttt{SQL Server (SQLEXPRESS)} está detenido o el
nombre de la instancia es otro &
Iniciar el servicio en \texttt{services.msc} y comprobar el nombre con SSMS
(paso M4). \\
\midrule
\texttt{docker compose up} falla con \emph{Bind for 0.0.0.0:8080 failed: port is already
allocated} (o 1433) &
Puerto ocupado por otro programa o por otra instancia de los contenedores &
Liberar el puerto (\texttt{netstat -ano | findstr :8080}) o cambiar el puerto
publicado en \texttt{docker-compose.yml}, por ejemplo \texttt{8081:8080}. \\
\midrule
\texttt{dotnet run} falla con \emph{address already in use} en el 5093 &
Otra instancia de la aplicación sigue en ejecución &
Cerrarla con \texttt{Ctrl+C} o terminar el proceso \texttt{personapi-dotnet}. \\
\midrule
\texttt{docker compose up} se detiene con \emph{Defina SA\_PASSWORD en .env} &
Falta el archivo \texttt{.env} o la variable está vacía &
Repetir el paso M12. \\
\midrule
\texttt{sqlserver} nunca llega a \emph{healthy} y su log rechaza la contraseña &
La contraseña no cumple la política de complejidad de SQL Server &
Poner una contraseña más fuerte en \texttt{.env} y ejecutar
\texttt{docker compose down -v} y luego \texttt{up --build -d}. \\
\midrule
\texttt{docker} responde \emph{cannot connect to the Docker daemon} &
Docker Desktop no está abierto o su motor no arrancó &
Abrir Docker Desktop y esperar a \emph{Engine running} (paso M6). \\
\bottomrule
\end{longtable}

Dos de estos fallos se reprodujeron a propósito en el equipo de verificación:

\begin{figure}[H]
\centering
\imagenOPendiente[0.85\textwidth]{figuras/manual/f-a-falta-sa-password.png}{\texttt{SA\_PASSWORD} vacía; Compose se detiene con el mensaje \emph{Defina SA\_PASSWORD en .env}}
\caption{Fallo: \texttt{SA\_PASSWORD} vacía; Compose se detiene con el mensaje \emph{Defina SA\_PASSWORD en .env}.}
\end{figure}
\begin{figure}[H]
\centering
\imagenOPendiente[0.85\textwidth]{figuras/manual/f-b-puerto-ocupado.png}{puerto ocupado; \texttt{dotnet run} en el 8080, que ya usa el contenedor, termina con \emph{address already in use}}
\caption{Fallo: puerto ocupado; \texttt{dotnet run} en el 8080, que ya usa el contenedor, termina con \emph{address already in use}.}
\end{figure}

\subsubsection{Registro de verificación}
El 27 de octubre de 2026 se ejecutaron los dos caminos en un equipo con Windows 11 Pro,
con las versiones de la Tabla~\ref{tab:manual-versiones}:

\begin{itemize}[itemsep=2pt]
  \item \textbf{Camino B desde cero.} Se eliminaron contenedores, red y volumen
        (\texttt{docker compose down -v}) y se ejecutó \texttt{docker compose up --build -d}
        sobre un clon nuevo del repositorio. \texttt{sqlserver} quedó \emph{healthy},
        \texttt{db-init} terminó con código 0 tras ejecutar los dos scripts y
        \texttt{personapi} escuchó en el 8080. Las comprobaciones de la
        Tabla~\ref{tab:manual-verificacion} dieron los resultados esperados, y la
        consulta a la base del contenedor mostró 3 filas por tabla. Una segunda ejecución
        de \texttt{docker compose up -d} no duplicó datos.
  \item \textbf{Camino A.} Con la base local creada por los scripts, \texttt{dotnet run}
        escuchó en \url{http://localhost:5093}; \texttt{/} y \texttt{/swagger}
        respondieron 200, y \texttt{/api/\allowbreak personas} y \texttt{/api/\allowbreak profesiones}
        devolvieron los datos de \texttt{persona\_db}. Sin cadena de conexión y con una
        base inexistente se reprodujeron los dos primeros fallos de la
        Tabla~\ref{tab:manual-fallos}.
\end{itemize}

\textbf{Capturas.} Las figuras de este manual son capturas reales de ese mismo equipo,
tomadas el 27 de octubre de 2026 al ejecutar cada paso: las ventanas de PowerShell,
SSMS 22 y Docker Desktop, y Google Chrome sobre la aplicación en los puertos 5093 y
8080. El 28 de octubre de 2026 se añadieron las de instalación con \texttt{winget},
las del Visual Studio Installer, el CRUD por Swagger, los recursos de Docker y dos
fallos reproducidos. En los pasos de instalación (M1 a M6), \texttt{winget install} se
ejecutó de verdad sobre herramientas que ya estaban instaladas, así que las capturas
muestran cómo detecta la versión existente y cancela la reinstalación; las pantallas
de Modificar del Visual Studio Installer se cerraron sin cambios, y el asistente gráfico
de Git se recorrió con los valores por defecto y se canceló antes de \emph{Install}. Los
asistentes de M2 a M6 (SDK de .NET 10.0.403, SQL Server 2025 Express, SSMS 22, \texttt{sqlcmd},
Visual Studio Community 2026 y Docker Desktop 4.97.0), descargados con \texttt{winget download},
también se recorrieron y se cerraron o cancelaron antes de instalar; las versiones de la
tabla de requisitos no cambiaron. Ninguna captura
muestra la contraseña de \texttt{sa}.

\textbf{Pendiente:} la ejecución de los pasos M1 a M6 en un equipo sin ninguna de las
herramientas instaladas. En la fecha del registro, los recursos de teléfonos y estudios y las vistas CRUD estaban en
desarrollo en otras historias; la verificación se limitó a lo integrado en ese momento.

\subsection{Pruebas de funcionamiento}

\begin{figure}[H]
\centering
\imagenOPendiente[\textwidth]{figuras/pruebas/f1-swagger.jpg}{Swagger UI con las 20 operaciones agrupadas por recurso}
\caption{Swagger UI (\texttt{/swagger}) con las 20 operaciones REST agrupadas por recurso:
estudios, personas, profesiones y teléfonos, cinco cada uno.}
\label{fig:prueba-swagger}
\end{figure}

\begin{figure}[H]
\centering
\imagenOPendiente[\textwidth]{figuras/pruebas/f2-post-personas-201.jpg}{\texttt{POST /api/personas} con respuesta \texttt{201 Created}}
\caption{Ejecución de \texttt{POST /api/personas} desde Swagger UI: respuesta \texttt{201} con
la persona creada y la cabecera \texttt{Location} hacia \texttt{/api/personas/1040404040}.}
\label{fig:prueba-post}
\end{figure}

\begin{figure}[H]
\centering
\imagenOPendiente[\textwidth]{figuras/pruebas/f3-personas-index.jpg}{Vista \texttt{/Personas} con los registros del script DML}
\caption{Vista \texttt{/Personas} (Index) con las tres personas que carga
\texttt{scripts/02\_dml.sql}.}
\label{fig:prueba-personas}
\end{figure}

\begin{figure}[H]
\centering
\imagenOPendiente[\textwidth]{figuras/pruebas/f4-telefono-create.jpg}{Formulario de creación de teléfono con el selector de personas}
\caption{Formulario \texttt{/Telefonos/Create}: el selector de dueño ofrece las personas
registradas (se muestra desplegado como lista para que se vean las opciones).}
\label{fig:prueba-telefono}
\end{figure}

\begin{figure}[H]
\centering
\imagenOPendiente[\textwidth]{figuras/pruebas/f5-docker-compose-ps.png}{Contenedores en ejecución (\texttt{docker compose ps})}
\caption{\texttt{docker compose ps}: \texttt{personapi} publica el puerto 8080 y
\texttt{sqlserver} el 1433 en estado \emph{healthy}; \texttt{db-init} ya terminó y no aparece.}
\label{fig:prueba-compose}
\end{figure}

\begin{table}[H]
\centering
\caption{Matriz de pruebas funcionales.}
\begin{tabularx}{\textwidth}{c X c c}
\toprule
\textbf{\#} & \textbf{Caso} & \textbf{Esperado} & \textbf{Resultado} \\
\midrule
1 & Listar personas (API y vista) & 200 / tabla & \checkmark \\
2 & Crear persona con cc nueva & 201 & \checkmark \\
3 & Crear persona con cc existente & 409 & \checkmark \\
4 & Crear persona con género distinto de M/F & 400 & \checkmark \\
5 & Actualizar persona existente & 204 & \checkmark \\
6 & Consultar persona inexistente & 404 & \checkmark \\
7 & Eliminar persona con teléfonos asociados & 409 & \checkmark \\
8 & Crear teléfono para dueño inexistente & 400 & \checkmark \\
9 & CRUD completo de profesiones & 200/201/204 & \checkmark \\
10 & CRUD de estudios con llave compuesta & 200/201/204 & \checkmark \\
11 & Despliegue completo con \texttt{docker compose up} & app en :8080 & \checkmark \\
\bottomrule
\end{tabularx}
\end{table}

\subsection{Paso 9 y 10. Publicación y etiqueta}
\begin{lstlisting}[language=bash,numbers=none]
git add .
git commit -m "Laboratorio 1: monolito MVC + DAO con .NET 10 y SQL Server 2025"
git push origin main
git tag -a v1.0.0 -m "Entrega Laboratorio 1"
git push origin v1.0.0
\end{lstlisting}

Repositorio: \url{\repoURL} \quad--\quad Tag: \texttt{\repoTag}

% ---------------------------------------------------------------------
%  CONCLUSIONES
% ---------------------------------------------------------------------
\section{Conclusiones y lecciones aprendidas}

\subsection{Conclusiones}
\begin{itemize}[itemsep=4pt]
  \item Un monolito bien estructurado en capas ofrece una separación de responsabilidades
        comparable a la de arquitecturas más complejas, con un costo operativo mucho menor:
        un solo artefacto, un solo despliegue y comunicación en proceso. Para un dominio
        pequeño como el de este laboratorio es la opción más adecuada.
  \item El patrón MVC permitió que una misma lógica de negocio sirviera dos interfaces
        distintas (vistas Razor y API REST) sin duplicar código: solo cambia el controlador
        y el formato de la respuesta.
  \item El patrón DAO aisló por completo a los controladores de Entity Framework. Gracias a
        la inyección de dependencias, cambiar la implementación de persistencia solo
        requeriría modificar una línea de registro en \texttt{Program.cs}.
  \item El enfoque \emph{Database First} con \texttt{Scaffold-DbContext} redujo
        drásticamente el código repetitivo: a partir del esquema se obtuvieron entidades,
        relaciones y contexto listos para usar.
  \item Docker Compose hizo reproducible el entorno completo (base de datos, scripts y
        aplicación) con un solo comando, independientemente del sistema operativo de cada
        integrante.
\end{itemize}

\subsection{Lecciones aprendidas}
\begin{itemize}[itemsep=4pt]
  \item \textbf{Migrar un esquema entre motores no es automático}: tipos como
        \texttt{ENUM} o \texttt{TEXT} de MySQL no existen o están obsoletos en SQL Server y
        deben reemplazarse por equivalentes (\texttt{CHECK}, \texttt{VARCHAR(MAX)}).
  \item \textbf{Las versiones importan}: .NET 10 exige Visual Studio 2026; la guía, basada
        en versiones anteriores, debió adaptarse. Elegir una versión LTS asegura soporte
        durante la vida del proyecto.
  \item \textbf{Las relaciones bidireccionales generan ciclos} al serializar a JSON
        (\texttt{Persona} $\rightarrow$ \texttt{Telefonos} $\rightarrow$ \texttt{Persona});
        se resolvió con \texttt{ReferenceHandler.IgnoreCycles}. Una alternativa más limpia
        es exponer DTOs en lugar de entidades.
  \item \textbf{La autenticación de Windows no funciona en contenedores Linux}: la cadena
        de conexión debe externalizarse y sobrescribirse por variables de entorno según el
        ambiente.
  \item \textbf{Las llaves foráneas sin cascada deben manejarse en la aplicación}: en
        lugar de dejar que la excepción de SQL llegue al usuario, se valida previamente y
        se responde con un código HTTP significativo (\texttt{409 Conflict}).
  \item \textbf{Las llaves compuestas complican el enrutamiento}: la entidad
        \texttt{estudios} obliga a diseñar rutas con dos parámetros y vistas que
        identifican el registro por ambos valores.
\end{itemize}

\subsection{Preguntas para el grupo}
\begin{enumerate}[itemsep=6pt]
  \item Entity Framework Core ya implementa internamente los patrones \emph{Repository}
        (\texttt{DbSet}) y \emph{Unit of Work} (\texttt{DbContext}). ¿Se justifica agregar
        una capa DAO propia encima de EF Core, o se trata de una abstracción redundante?
        ¿En qué escenarios aporta valor real (pruebas, cambio de motor, consultas
        complejas) y en cuáles solo añade código?
  \item Si el sistema creciera (por ejemplo, millones de personas y equipos de desarrollo
        independientes para teléfonos y estudios), ¿qué señales indicarían que es momento de
        migrar este monolito a microservicios, y qué decisiones del diseño actual
        facilitarían o dificultarían esa transición?
\end{enumerate}

% ---------------------------------------------------------------------
%  REFERENCIAS
% ---------------------------------------------------------------------
\begin{thebibliography}{99}
\addcontentsline{toc}{section}{Referencias}

\bibitem{bass} L. Bass, P. Clements y R. Kazman, \emph{Software Architecture in Practice},
  4.\textsuperscript{a} ed. Boston: Addison-Wesley, 2021.

\bibitem{richardson} C. Richardson, \emph{Microservices Patterns: With Examples in Java}.
  Shelter Island: Manning, 2018. Patrón \enquote{Monolithic Architecture}:
  \url{https://microservices.io/patterns/monolithic.html}

\bibitem{msarchguide} Microsoft Patterns \& Practices, \emph{Application Architecture Guide
  2.0: Designing Applications on the .NET Platform}. Microsoft, 2009.

\bibitem{reenskaug} T. Reenskaug, \enquote{Models -- Views -- Controllers}, Xerox PARC,
  Technical Note, 1979. \url{https://folk.universitetetioslo.no/trygver/themes/mvc/mvc-index.html}

\bibitem{aspnetmvc} Microsoft, \enquote{Overview of ASP.NET Core MVC}, Microsoft Learn.
  \url{https://learn.microsoft.com/aspnet/core/mvc/overview}

\bibitem{corej2ee} D. Alur, J. Crupi y D. Malks, \emph{Core J2EE Patterns: Best Practices
  and Design Strategies}, 2.\textsuperscript{a} ed. Prentice Hall, 2003. Patrón
  \enquote{Data Access Object}.

\bibitem{fowler} M. Fowler, \emph{Patterns of Enterprise Application Architecture}.
  Boston: Addison-Wesley, 2002.

\bibitem{dotnetpolicy} Microsoft, \enquote{Directiva de soporte técnico de .NET y .NET
  Core}. \url{https://dotnet.microsoft.com/es-es/platform/support/policy/dotnet-core}

\bibitem{efcore} Microsoft, \enquote{Entity Framework Core} y \enquote{Scaffolding
  (Reverse Engineering)}, Microsoft Learn.
  \url{https://learn.microsoft.com/ef/core/managing-schemas/scaffolding/}

\bibitem{sqlserver} Microsoft, \enquote{What's new in SQL Server 2025}, Microsoft Learn.
  \url{https://learn.microsoft.com/sql/sql-server/what-s-new-in-sql-server-2025}

\bibitem{fielding} R. T. Fielding, \emph{Architectural Styles and the Design of
  Network-based Software Architectures}, tesis doctoral, University of California, Irvine,
  2000.

\bibitem{openapi} OpenAPI Initiative, \enquote{OpenAPI Specification v3.1.0}.
  \url{https://spec.openapis.org/oas/v3.1.0}

\bibitem{swashbuckle} Microsoft, \enquote{Get started with Swashbuckle and ASP.NET Core},
  Microsoft Learn. \url{https://learn.microsoft.com/aspnet/core/tutorials/getting-started-with-swashbuckle}

\bibitem{docker} Docker Inc., \enquote{Docker overview} y \enquote{Docker Compose
  overview}. \url{https://docs.docker.com/get-started/docker-overview/}

\bibitem{mssqldocker} Microsoft, \enquote{Quickstart: Run SQL Server Linux container
  images with Docker}. \url{https://learn.microsoft.com/sql/linux/quickstart-install-connect-docker}

\bibitem{deitel} P. Deitel y H. Deitel, \emph{Programación en la web}. Pearson Educación,
  2008.

\bibitem{hdeleon1} hdeleon.net, \enquote{Video tutorial ASP.NET Core MVC con Entity
  Framework}. \url{https://www.youtube.com/watch?v=6nT-RjMEG0o}

\bibitem{hdeleon2} hdeleon.net, \enquote{Video tutorial Web API con .NET}.
  \url{https://www.youtube.com/watch?v=28LjewDjaz4}

\bibitem{guamann} V. Guamán, \enquote{API con ASP.NET MVC 6 y SQL Server mediante Entity
  Framework Core 6 -- Code First}, partes 1 y 2, DEV Community.
  \url{https://dev.to/veronicaguamann}

\bibitem{csharpcorner} C\# Corner, \enquote{Building ASP.NET Web API in .NET Core with
  Entity Framework}.
  \url{https://www.c-sharpcorner.com/article/building-asp-net-web-api-in-net-core-with-entity-framework/}

\bibitem{curso} Material del curso \asignatura: \enquote{8.1 Introducción a la Plataforma
  .NET} y \enquote{8.2 Laboratorio .NET}.

\end{thebibliography}

\end{document}